\documentclass[11pt,a4paper]{article}
% =====================================================================
%  Gemeinsame Vorlage der Arbeitsblätter – 1. Formale Sprachen (Q13)
%  (gleiches Layout wie informatik/13/material/grundwissen/ab-vorlage.tex
%  und informatik/12/rekursive-datenstrukturen/arbeitsblaetter/ab-vorlage.tex)
%  Einbinden in jedem Blatt direkt nach \documentclass: % =====================================================================
%  Gemeinsame Vorlage der Arbeitsblätter – 1. Formale Sprachen (Q13)
%  (gleiches Layout wie informatik/13/material/grundwissen/ab-vorlage.tex
%  und informatik/12/rekursive-datenstrukturen/arbeitsblaetter/ab-vorlage.tex)
%  Einbinden in jedem Blatt direkt nach \documentclass: % =====================================================================
%  Gemeinsame Vorlage der Arbeitsblätter – 1. Formale Sprachen (Q13)
%  (gleiches Layout wie informatik/13/material/grundwissen/ab-vorlage.tex
%  und informatik/12/rekursive-datenstrukturen/arbeitsblaetter/ab-vorlage.tex)
%  Einbinden in jedem Blatt direkt nach \documentclass: \input{ab-vorlage}
%  Übersetzen mit XeLaTeX, LuaLaTeX, pdfLaTeX oder Tectonic.
%
%  Blatt:   tectonic ab-1-2-formale-sprachen.tex
%  Lösung:  tectonic ab-1-2-formale-sprachen-lsg.tex
%           (Datei enthält nur: \def\mitloesung{}\input{ab-1-2-formale-sprachen};
%            füllt alle Lücken rot aus; siehe \lsg, \lsglinien, \lsgoder, lsgbild, \karolsg)
% =====================================================================
\usepackage[a4paper,left=16mm,right=16mm,top=23mm,bottom=13mm,headheight=12mm,headsep=3mm,footskip=6mm]{geometry}
\usepackage{iftex}
\ifPDFTeX
  \usepackage[T1]{fontenc}
  \usepackage[utf8]{inputenc}
  \usepackage{helvet}
  \usepackage[scaled=0.86]{beramono}
  \renewcommand{\familydefault}{\sfdefault}
\else
  \usepackage{fontspec}
  \setmainfont{texgyreheros}[Extension=.otf,UprightFont=*-regular,BoldFont=*-bold,ItalicFont=*-italic,BoldItalicFont=*-bolditalic]
  \setmonofont{DejaVuSansMono}[Extension=.ttf,UprightFont=*,BoldFont=*-Bold,ItalicFont=*-Oblique,BoldItalicFont=*-BoldOblique,Scale=0.86]
\fi
\usepackage[ngerman,shorthands=off]{babel}
\usepackage{amssymb}
\usepackage{xcolor,graphicx,tabularx,array,enumitem,fancyhdr,tikz,multicol,calc,etoolbox,pifont}
\usepackage[most]{tcolorbox}
\usetikzlibrary{arrows.meta,positioning,calc}
\usepackage[hidelinks]{hyperref}
\setlength{\parindent}{0pt}
\setlength{\parskip}{3pt}
\setlist{nosep,leftmargin=*}

% ---------- Lösungsschalter ----------
\newif\ifloesung
\ifdefined\mitloesung\loesungtrue\fi
\newcommand{\lsgfarbe}{\color{red!75!black}}

% ---------- Kopf- und Fußzeile (einheitliche Form) ----------
\newcommand{\lehrkraft}{Hau}
\newcommand{\themenbereich}{1. Formale Sprachen}
\newcommand{\abnummer}{}
\pagestyle{fancy}\fancyhf{}
\renewcommand{\headrulewidth}{0.4pt}
\fancyhead[L]{\small Klasse 13\\Datum:}
\fancyhead[C]{\small Themenbereich:\\\themenbereich}
\fancyhead[R]{\small \lehrkraft\ifloesung\\{\lsgfarbe\bfseries Lösung}\fi}
\fancyfoot[C]{\scriptsize Informatik 13 gA\,·\,Blatt \abnummer\ifloesung\ (Lösung)\fi\,·\,Seite \thepage}
\newcommand{\abtitel}[2]{\renewcommand{\abnummer}{#1}\begin{center}\Large\bfseries #1\quad\underline{#2}\end{center}\vspace{-2mm}}

% ---------- Boxen ----------
\newenvironment{aufgabe}[2][]{\begin{tcolorbox}[enhanced,breakable,colback=white,colframe=black,boxrule=0.7pt,arc=3mm,left=2mm,right=2mm,top=1mm,bottom=1.2mm,before skip=2mm,after skip=2mm]\textbf{Aufgabe #2}\ifstrempty{#1}{}{ \textit{#1}}\par\vspace{0.3mm}}{\end{tcolorbox}}
\newtcolorbox{merke}[1][Merke]{enhanced,breakable,colback=green!5,colframe=green!40!black,boxrule=0.8pt,arc=2mm,left=2mm,right=2mm,top=1mm,bottom=1.2mm,before skip=2mm,after skip=2mm,title={#1},fonttitle=\bfseries,coltitle=white,toptitle=0.3mm,bottomtitle=0.3mm}
\newtcolorbox{hinweis}{enhanced,breakable,colback=yellow!12,colframe=orange!70!black,boxrule=0.6pt,arc=2mm,left=2mm,right=2mm,top=1mm,bottom=1mm,before skip=2mm,after skip=2mm}

% ---------- Lücken, Linien, Karos ----------
\newcommand{\luecke}[1][5.25cm]{\rule[-0.5ex]{#1}{0.4pt}}
\newcommand{\linien}[1]{\par\vspace{1.5mm}\foreach \i in {1,...,#1}{\noindent\rule[0pt]{\linewidth}{0.3pt}\par\vspace{5.5mm}}\vspace{-3mm}}
\newcommand{\code}[1]{\texttt{#1}}
\newcommand{\eps}{$\varepsilon$}

% Lücke mit Lösung: \lsg[Mindestbreite]{Lösungstext}
%  Die Lösung steht (ohne eigene Höhe/Tiefe) auf derselben Linie wie im Blatt: Das Layout verschiebt sich nicht,
%  solange der Text nicht breiter als die Mindestbreite ist.
\newlength{\lsgbreite}
\newcommand{\lsg}[2][5.25cm]{\leavevmode\ifloesung\settowidth{\lsgbreite}{\,#2\,}\ifdim\lsgbreite<#1\setlength{\lsgbreite}{#1}\fi\rlap{\smash{\raisebox{0.8pt}{\makebox[\lsgbreite][l]{\lsgfarbe\,#2}}}}\luecke[\lsgbreite]\else\luecke[#1]\fi}
% Lücke bis zum Zeilenende (beendet den Absatz): \lsgrest{Lösungstext} bzw. \restzeile
\newcommand{\lsgrest}[1]{\leavevmode\unskip\hspace{1.5mm}\ifloesung\rlap{\smash{\raisebox{0.8pt}{\lsgfarbe\,#1}}}\fi\leaders\hbox{\rule[-0.5ex]{0.5pt}{0.4pt}}\hfill\null\par}
\newcommand{\restzeile}{\lsgrest{}}
% Schreiblinien mit Lösung: \lsglinien{Anzahl}{Lösungstext} (gleiche Höhe wie \linien:
%  je Linie \baselineskip + 5,5 mm + \parskip, ausgemessen)
\newlength{\lsglinienhoehe}
\newcommand{\lsglinien}[2]{\ifloesung\par\vspace{1.5mm}\setlength{\lsglinienhoehe}{\dimexpr #1\dimexpr\baselineskip+5.5mm+\parskip\relax-16.25pt\relax}\noindent\begin{minipage}[t][\lsglinienhoehe][t]{\linewidth}\lsgfarbe #2\end{minipage}\par\vspace{-1.5mm}\else\linien{#1}\fi}
% Beliebiger Inhalt: \lsgoder{nur in der Lösung}{nur im Blatt}
\newcommand{\lsgoder}[2]{\ifloesung\leavevmode#1\else#2\fi}
% Text nur in der Lösung sichtbar, belegt im Blatt aber denselben Platz
\newcommand{\lsgtext}[1]{\leavevmode\ifloesung{\lsgfarbe#1}\else\phantom{#1}\fi}
% Tabellenzelle zum Hineinschreiben (8 mm hoch, ohne Schreiblinie; Lösung mit \lsgtext{…})
\newcommand{\schreibzeile}{\rule[-2.8mm]{0pt}{8mm}}
% Ankreuzfeld: \kreuz{1} = angekreuzt (nur in der Lösung), \kreuz{0} = leer
\newcommand{\kreuz}[1]{\leavevmode\ifloesung\ifnum#1=1 {\lsgfarbe$\boxtimes$}\else$\square$\fi\else$\square$\fi}

% ---------- Syntaxdiagramme (Maße in mm) ----------
%  Terminal rund (orange), Nichtterminal eckig (grün), Linien mit Pfeilen:
%  \draw[sdl,->] (0,0) -- ++(6,0) node[nt,anchor=west] (a) {Augen};   (Pfeil endet am Kasten)
\tikzset{
  sd/.style={x=1mm,y=1mm,>={Stealth[length=1.6mm,width=1.3mm]},font=\small},
  sdl/.style={line width=0.55pt,rounded corners=1.2mm},
  te/.style={draw,rounded corners=2.1mm,minimum height=4.2mm,minimum width=4.2mm,inner sep=1pt,font=\ttfamily\small,fill=orange!18},
  nt/.style={draw,sharp corners,minimum height=4.2mm,inner sep=1.6pt,font=\small,fill=green!12},   % sharp corners: sonst erbt der Kasten die runden Ecken der Linie (sdl)
  sdname/.style={anchor=south west,font=\small\bfseries,inner sep=0pt},
  punkte/.style={font=\small,inner sep=0pt},
  % Ableitungsbaum
  baum/.style={font=\small,edge from parent/.style={draw,line width=0.5pt,-{Stealth[length=1.4mm]}},level distance=8.5mm},
}
% Lösungszeichnung: im Blatt unsichtbar, belegt aber denselben Platz
\ifloesung
  \tikzset{lsgbild/.style={red!75!black,te/.append style={fill=red!5},nt/.append style={fill=red!5}}}
\else
  \tikzset{lsgbild/.style={draw opacity=0,fill opacity=0,text opacity=0}}
\fi
% Karofläche, in der Lösung mit roter Zeichnung: \karolsg{Anzahl Karozeilen}{TikZ-Code}
% Koordinaten der Zeichnung in mm, Ursprung links oben, y nach unten negativ.
\newcommand{\karolsg}[2]{\par\vspace{1mm}\noindent\begin{tikzpicture}
  \draw[step=5mm,gray!45,thin] (0,0) grid (\linewidth-1pt,#1*5mm);
  \useasboundingbox (0,0) rectangle (\linewidth-1pt,#1*5mm);
  \ifloesung\begin{scope}[sd,shift={(0,#1*5mm)},lsgbild]#2\end{scope}\fi
\end{tikzpicture}\par}

% ---------- Bausteine für Syntaxdiagramme (innerhalb tikzpicture[sd]) ----------
%  Jeder Baustein beginnt an einem Punkt (Knoten-Anker oder Koordinate) und
%  legt die Koordinate <name>-o an, an der es weitergeht.
% Aneinanderreihung: \sdnext{name}{von}{te|nt}{Text}   (weiter bei name.east)
\newcommand{\sdnext}[4]{\draw[sdl,->] (#2) -- ++(4,0) node[#3,anchor=west] (#1) {#4};\coordinate (#1-o) at (#1.east);}
% Option: \sdopt{name}{von}{te|nt}{Text}
\newcommand{\sdopt}[4]{\path (#2) ++(4,-5.5) node[#3,anchor=west] (#1) {#4};
  \coordinate (#1-x) at ($(#1.east)+(3.5,0)$); \coordinate (#1-o) at (#1-x |- #2);
  \draw[sdl,->] (#2) |- (#1.west); \draw[sdl,->] (#1.east) -| (#1-o); \draw[sdl] (#2) -- (#1-o);}
% Wiederholung, mindestens einmal: \sdwdh{name}{von}{te|nt}{Text}
\newcommand{\sdwdh}[4]{\draw[sdl,->] (#2) -- ++(6,0) node[#3,anchor=west] (#1) {#4};
  \coordinate (#1-o) at ($(#1.east)+(4,0)$);
  \draw[sdl] (#1.east) -- (#1-o); \draw[sdl,->] (#1-o) -- ++(0,-5.5) -| ($(#2)+(2.5,0)$);}
% Wiederholung, beliebig oft (auch keinmal): \sdwdhn{name}{von}{te|nt}{Text}
\newcommand{\sdwdhn}[4]{\path (#2) ++(6,-5.5) node[#3,anchor=west] (#1) {#4};
  \coordinate (#1-x) at ($(#1.east)+(4,0)$); \coordinate (#1-o) at (#1-x |- #2);
  \draw[sdl] (#2) -- (#1-o); \draw[sdl,->] (#1-o) |- (#1.east); \draw[sdl,->] (#1.west) -| ($(#2)+(2.5,0)$);}
% Alternative: \sdalt{name}{von}{Breite bis zum Zusammenführen}{Zeilenabstand}{te/a, nt/Ziffer, pk/, ...}
%  pk/ setzt senkrechte Pünktchen (…) statt eines Zweigs
\newcommand{\sdalt}[5]{\coordinate (#1-o) at ($(#2)+(#3,0)$);
  \foreach \typ/\txt [count=\i from 0] in {#5}{%
    \ifdefstring{\typ}{pk}{\node[punkte] at ($(#2)+(6.1,-\i*#4+0.4)$) {$\vdots$};}{%
      \ifdefstring{\typ}{te}{\tikzset{sdakt/.style={te}}}{\tikzset{sdakt/.style={nt}}}%
      \node[sdakt,anchor=west] (#1\i) at ($(#2)+(4,-\i*#4)$) {\txt};
      \ifnum\i=0\relax \draw[sdl,->] (#2) -- (#1\i.west); \draw[sdl] (#1\i.east) -- (#1-o);
      \else \draw[sdl,->] (#2) |- (#1\i.west); \draw[sdl,->] (#1\i.east) -| (#1-o);\fi}}}
% Ende eines Diagramms
\newcommand{\sdende}[1]{\draw[sdl,->] (#1) -- ++(5,0);}

% ---------- Die Smiley-Grammatik (Blatt 1.3 und 1.4) ----------
\newcommand{\sdSmileyfolge}{\begin{tikzpicture}[sd,baseline=0mm]
  \node[sdname] at (0,3.3) {Smileyfolge}; \coordinate (s) at (0,0);
  \sdwdh{sm}{s}{nt}{Smiley}\sdende{sm-o}\end{tikzpicture}}
\newcommand{\sdSmiley}{\begin{tikzpicture}[sd,baseline=0mm]
  \node[sdname] at (0,3.3) {Smiley}; \coordinate (s) at (0,0);
  \sdnext{au}{s}{nt}{Augen}\sdnext{mi}{au-o}{nt}{Mitte}\sdnext{mu}{mi-o}{nt}{Mund}\sdende{mu-o}\end{tikzpicture}}
\newcommand{\sdAugen}{\begin{tikzpicture}[sd,baseline=0mm]
  \node[sdname] at (0,3.3) {Augen}; \coordinate (s) at (0,0);
  \sdalt{a}{s}{12}{5.5}{te/:, te/;}\sdende{a-o}\end{tikzpicture}}
\newcommand{\sdMitte}{\begin{tikzpicture}[sd,baseline=0mm]
  \node[sdname] at (0,3.3) {Mitte}; \coordinate (s) at (0,0);
  \sdopt{n}{s}{te}{-}\sdende{n-o}\end{tikzpicture}}
\newcommand{\sdMund}{\begin{tikzpicture}[sd,baseline=0mm]
  \node[sdname] at (0,3.3) {Mund}; \coordinate (s) at (0,0);
  \sdalt{m}{s}{12}{5.5}{te/), te/{(}, te/D, te/P, te/*}\sdende{m-o}\end{tikzpicture}}

% ---------- Ergänzungen für Programmier-Blätter (ab 1.6; aus informatik/10/java/arbeitsblaetter/ab-vorlage.tex) ----------
%  Nur angehängt: Die Makros oben bleiben unverändert, damit die PDFs von 1.1 bis 1.5 gleich bleiben.
%  \lueckentext   am Anfang eines Absatzes/Kastens mit Lücken: Zeilenabstand 1,5, Flattersatz
%  \nr{1}         Schrittnummer ① wie auf der Website und im Code
%  java           Listing:  \begin{java}[numbers=left] … \end{java}
%  javaluecke     Lückencode mit größerem Zeilenabstand, Lücken als (*\lsg[6.75cm]{Lösung}*)
\newcommand{\lueckentext}{\linespread{1.5}\selectfont\raggedright}
\newcommand{\nr}[1]{\ding{\numexpr171+#1\relax}}
\usepackage{listings}
\lstset{language=Java,basicstyle=\ttfamily\small,keywordstyle=\bfseries,commentstyle=\color{gray!80!black},stringstyle=\color{black},showstringspaces=false,columns=flexible,keepspaces=true,tabsize=3,frame=single,framerule=0.4pt,rulecolor=\color{black},xleftmargin=2mm,framexleftmargin=1mm,aboveskip=1.5mm,belowskip=1.5mm,breaklines=true,numberstyle=\tiny\color{gray},numbersep=3mm,escapeinside={(*}{*)},
 literate={ä}{{\"a}}1 {ö}{{\"o}}1 {ü}{{\"u}}1 {Ä}{{\"A}}1 {Ö}{{\"O}}1 {Ü}{{\"U}}1 {ß}{{\ss}}1 {–}{{--}}1 {…}{{\dots}}1}
\lstnewenvironment{java}[1][]{\lstset{#1}}{}
\lstnewenvironment{javaluecke}[1][]{\lstset{basicstyle=\linespread{1.5}\small\ttfamily,#1}}{}

%  Übersetzen mit XeLaTeX, LuaLaTeX, pdfLaTeX oder Tectonic.
%
%  Blatt:   tectonic ab-1-2-formale-sprachen.tex
%  Lösung:  tectonic ab-1-2-formale-sprachen-lsg.tex
%           (Datei enthält nur: \def\mitloesung{}\documentclass[11pt,a4paper]{article}
\input{ab-vorlage}
% Blatt:   tectonic ab-1-2-formale-sprachen.tex
% Lösung:  tectonic ab-1-2-formale-sprachen-lsg.tex
% Übungsfassung: Die Begriffe (Alphabet, Wort, Länge, ε, A*, L ⊆ A*) führt Blatt 1.1 ein, hier wird nur geübt.
% Progression: Begriffe sicher anwenden (ε!) -> Wörter zählen |A|^n, endliche Sprache in unendlichem A*
%              -> Sprachen durch Eigenschaften festlegen -> Alltag (Syntax, nicht Semantik)
%              -> Smileys in Worten beschreiben: mehrdeutig -> Regeln nötig (Brücke zu 1.3)
\newcommand{\zeile}{\rule[-2mm]{0pt}{6.6mm}}
\newcommand{\hoch}{\rule[-1.6mm]{0pt}{6mm}}
\begin{document}
\abtitel{1.2}{Formale Sprachen -- Übung}

Die Begriffe kennst du aus 1.1 (Merkkasten \emph{Formale Sprachen}). Auf diesem Blatt übst du sie mit drei Alphabeten:\\[1mm]
\centerline{$A_1 = \{\code{0}, \code{1}\}$\qquad $A_2 = \{\code{a}, \code{b}, \code{c}, \dots, \code{z}\}$\qquad $A_3 = \{\code{:}, \code{;}, \code{-}, \code{)}, \code{(}, \code{D}, \code{P}, \code{*}\}$}

\begin{aufgabe}[Kurz geprüft]{1}
\textbf{a)} Kreuze alle richtigen Aussagen an.\\[1mm]
\renewcommand{\arraystretch}{1.25}
\begin{tabularx}{\linewidth}{@{}X@{\quad}X@{}}
\kreuz{1} $\code{0} \in A_1$ & \kreuz{0} $\varepsilon \in A_1$\\
\kreuz{0} \code{ab} ist ein Wort über $A_1$ & \kreuz{1} $\varepsilon \in A_1^*$\\
\kreuz{0} $A_1^*$ ist endlich, weil $A_1$ endlich ist. & \kreuz{1} \code{:-)} ist ein Wort der Länge 3 über $A_3$.\\
\kreuz{1} Jedes Zeichen aus $A$ ist auch ein Wort über $A$. & \kreuz{0} Das leere Wort $\varepsilon$ ist ein Leerzeichen.\\
\end{tabularx}\\[1mm]
\textbf{b)} Es gibt genau ein Wort, das über \emph{jedem} Alphabet gebildet werden kann. Welches? Warum?\\[0.6mm]
\lsg[\linewidth]{Das leere Wort \eps: Es enthält kein Zeichen, also auch keines, das im Alphabet fehlen könnte.}
\end{aufgabe}

\begin{aufgabe}[Wie viele Wörter gibt es?]{2}
\textbf{a)} Schreibe alle Wörter über $A_1$ der Länge 0 bis 3 auf und zähle. Zähle auch für $A_3$ und allgemein.\\[1mm]
\renewcommand{\arraystretch}{1.1}
\begin{tabularx}{\linewidth}{|p{4.4cm}|p{0.8cm}|p{1.3cm}|p{2.6cm}|X|p{1.1cm}|}\hline
\textbf{Länge $n$} & \textbf{0} & \textbf{1} & \textbf{2} & \textbf{3} & \textbf{$n$}\\\hline
Wörter über $A_1$\rule[-6.5mm]{0pt}{11mm} & \lsgoder{\lsgfarbe \eps}{} & \lsgoder{\lsgfarbe \code{0}, \code{1}}{} & \lsgoder{\lsgfarbe \code{00}, \code{01}, \code{10}, \code{11}}{} & \lsgoder{\lsgfarbe \code{000}, \code{001}, \code{010}, \code{011}, \code{100}, \code{101}, \code{110}, \code{111}}{} & \\\hline
Anzahl über $A_1$\hoch & \lsgoder{\lsgfarbe 1}{} & \lsgoder{\lsgfarbe 2}{} & \lsgoder{\lsgfarbe 4}{} & \lsgoder{\lsgfarbe 8}{} & \lsgoder{\lsgfarbe $2^n$}{}\\\hline
Anzahl über $A_3$\hoch & \lsgoder{\lsgfarbe 1}{} & \lsgoder{\lsgfarbe 8}{} & \lsgoder{\lsgfarbe 64}{} & \lsgoder{\lsgfarbe 512}{} & \lsgoder{\lsgfarbe $8^n$}{}\\\hline
Anzahl bei $k$ Zeichen in $A$\hoch & \lsgoder{\lsgfarbe 1}{} & \lsgoder{\lsgfarbe $k$}{} & \lsgoder{\lsgfarbe $k^2$}{} & \lsgoder{\lsgfarbe $k^3$}{} & \lsgoder{\lsgfarbe $k^n$}{}\\\hline
\end{tabularx}\\[1.2mm]
\textbf{b)} Wie viele Farbcodes wie \code{\#9900BB} gibt es (1.1, Aufgabe 6)? Ist diese Sprache endlich oder unendlich?\\[0.6mm]
\lsg[\linewidth]{$16^6 = 16\,777\,216$ Farbcodes: Die Sprache ist endlich, obwohl $A_\text{Hex}^*$ unendlich ist.}
\end{aufgabe}

\begin{aufgabe}[Sprachen über $A_2$]{3}
\begin{minipage}[t]{0.585\linewidth}
\vspace{0pt}
$L_1$ = alle Wörter aus $A_2^*$ mit Länge kleiner gleich 4\\
$L_2$ = alle Wörter aus $A_2^*$, die mit \code{b} beginnen\\
$L_3$ = alle Wörter aus $A_2^*$, die das Teilwort \code{inf} enthalten\\[1mm]
\textbf{a)} Kreuze an, zu welcher Sprache das Wort gehört.\\[1mm]
\textbf{b)} Welche der drei Sprachen sind endlich? \lsg[2.7cm]{nur $L_1$}\\[1mm]
\textbf{c)} Finde ein Wort, das zu $L_1$, $L_2$ \emph{und} $L_3$ gehört:\\[0.6mm] \lsg[4.8cm]{\code{binf}}\\[1.2mm]
Gibt es mehrere? \lsg[3.9cm]{nein, nur \code{binf}}
\end{minipage}\hfill
\begin{minipage}[t]{0.39\linewidth}
\vspace{0pt}
\renewcommand{\arraystretch}{1.12}
\begin{tabularx}{\linewidth}{|X|c|c|c|}\hline
\textbf{Wort} & $L_1$ & $L_2$ & $L_3$\\\hline
\code{baum} & \kreuz{1} & \kreuz{1} & \kreuz{0}\\\hline
\code{informatik} & \kreuz{0} & \kreuz{0} & \kreuz{1}\\\hline
\eps & \kreuz{1} & \kreuz{0} & \kreuz{0}\\\hline
\code{inf} & \kreuz{1} & \kreuz{0} & \kreuz{1}\\\hline
\code{binfo} & \kreuz{0} & \kreuz{1} & \kreuz{1}\\\hline
\end{tabularx}
\end{minipage}
\end{aufgabe}

\newpage
\begin{aufgabe}[Formale Sprachen im Alltag]{4}
\textbf{a)} Ergänze die Tabelle.\\[1mm]
\renewcommand{\arraystretch}{1.1}
\begin{tabularx}{\linewidth}{|p{5.3cm}|X|p{2.4cm}|p{4.5cm}|}\hline
\textbf{Sprache} & \textbf{Alphabet} & \textbf{Wort aus $L$} & \textbf{Wort aus $A^*$, nicht aus $L$}\\\hline
Postleitzahlen\zeile & \lsgoder{\lsgfarbe \{0, 1, \dots, 9\}}{} & \code{83278} & \lsgoder{\lsgfarbe \code{824377} (6 Ziffern)}{}\\\hline
Binärzahlen ohne führende Null\zeile & \lsgoder{\lsgfarbe \{0, 1\}}{} & \lsgoder{\lsgfarbe \code{1011}}{} & \lsgoder{\lsgfarbe \code{0011}}{}\\\hline
Datum (TT.MM.JJJJ)\zeile & \lsgoder{\lsgfarbe \{0, 1, \dots, 9, .\}}{} & \lsgoder{\lsgfarbe \code{29.09.2026}}{} & \lsgoder{\lsgfarbe \code{29..092026}}{}\\\hline
\end{tabularx}\\[1.2mm]
\textbf{b)} Die Postleitzahl \code{99999} ist nicht vergeben. Gehört sie trotzdem zur Sprache „alle Wörter aus fünf Ziffern“? Was folgt daraus für Syntax und Semantik?
\lsglinien{2}{Ja, \code{99999} ist syntaktisch korrekt. Eine formale Sprache legt nur die \textbf{Syntax} fest; ob das Wort auch etwas bedeutet (\textbf{Semantik}), prüft sie nicht.}
\end{aufgabe}

\begin{aufgabe}[Die Sprache der Smileys]{5}
Im Chat werden Smileys aus den Zeichen von $A_3$ gebildet. Sie bilden eine Sprache $L_\text{Smiley} \subseteq A_3^*$.\\[1mm]
\textbf{a)} Kreuze an, welche Wörter du als Smiley erkennst.\\[1mm]
\kreuz{1} \code{:-)}\qquad \kreuz{1} \code{;D}\qquad \kreuz{1} \code{:-P}\qquad \kreuz{0} \code{)-:}\qquad \kreuz{0} \code{:-}\qquad \kreuz{1} \code{;*}\qquad \kreuz{0} \code{:--)}\qquad \kreuz{0} \code{DP}\qquad \kreuz{1} \code{:(}\\[1mm]
\textbf{b)} Beschreibe in Worten möglichst genau, welche Wörter aus $A_3^*$ zu $L_\text{Smiley}$ gehören.
\lsglinien{3}{Ein Smiley beginnt mit Augen (\code{:} oder \code{;}). Danach kann eine Nase \code{-} folgen, höchstens eine. Am Ende steht genau ein Mund: \code{)}, \code{(}, \code{D}, \code{P} oder \code{*}. Sonst folgt nichts.}
\textbf{c)} Vergleiche mit deinem Nachbarn: Gibt es Wörter, bei denen eure Beschreibungen verschieden entscheiden oder gar nicht entscheiden? Notiere ein Beispiel und woran es liegt.
\lsglinien{2}{z.\,B. \code{)-:} (gespiegelt?), \code{:--)} (wie viele Nasen?), \code{:-):-)} (mehrere Smileys?). Umgangssprache ist ungenau und mehrdeutig; man übersieht leicht Fälle.}
\end{aufgabe}

\begin{merke}[Erkenntnis]
Eine Beschreibung in Umgangssprache legt eine Sprache oft \lsg[4.2cm]{nicht eindeutig} fest. Nötig sind genaue \textbf{Regeln}, nach denen man für jedes Wort entscheiden kann, ob es zur Sprache gehört.
\end{merke}
\end{document}
;
%            füllt alle Lücken rot aus; siehe \lsg, \lsglinien, \lsgoder, lsgbild, \karolsg)
% =====================================================================
\usepackage[a4paper,left=16mm,right=16mm,top=23mm,bottom=13mm,headheight=12mm,headsep=3mm,footskip=6mm]{geometry}
\usepackage{iftex}
\ifPDFTeX
  \usepackage[T1]{fontenc}
  \usepackage[utf8]{inputenc}
  \usepackage{helvet}
  \usepackage[scaled=0.86]{beramono}
  \renewcommand{\familydefault}{\sfdefault}
\else
  \usepackage{fontspec}
  \setmainfont{texgyreheros}[Extension=.otf,UprightFont=*-regular,BoldFont=*-bold,ItalicFont=*-italic,BoldItalicFont=*-bolditalic]
  \setmonofont{DejaVuSansMono}[Extension=.ttf,UprightFont=*,BoldFont=*-Bold,ItalicFont=*-Oblique,BoldItalicFont=*-BoldOblique,Scale=0.86]
\fi
\usepackage[ngerman,shorthands=off]{babel}
\usepackage{amssymb}
\usepackage{xcolor,graphicx,tabularx,array,enumitem,fancyhdr,tikz,multicol,calc,etoolbox,pifont}
\usepackage[most]{tcolorbox}
\usetikzlibrary{arrows.meta,positioning,calc}
\usepackage[hidelinks]{hyperref}
\setlength{\parindent}{0pt}
\setlength{\parskip}{3pt}
\setlist{nosep,leftmargin=*}

% ---------- Lösungsschalter ----------
\newif\ifloesung
\ifdefined\mitloesung\loesungtrue\fi
\newcommand{\lsgfarbe}{\color{red!75!black}}

% ---------- Kopf- und Fußzeile (einheitliche Form) ----------
\newcommand{\lehrkraft}{Hau}
\newcommand{\themenbereich}{1. Formale Sprachen}
\newcommand{\abnummer}{}
\pagestyle{fancy}\fancyhf{}
\renewcommand{\headrulewidth}{0.4pt}
\fancyhead[L]{\small Klasse 13\\Datum:}
\fancyhead[C]{\small Themenbereich:\\\themenbereich}
\fancyhead[R]{\small \lehrkraft\ifloesung\\{\lsgfarbe\bfseries Lösung}\fi}
\fancyfoot[C]{\scriptsize Informatik 13 gA\,·\,Blatt \abnummer\ifloesung\ (Lösung)\fi\,·\,Seite \thepage}
\newcommand{\abtitel}[2]{\renewcommand{\abnummer}{#1}\begin{center}\Large\bfseries #1\quad\underline{#2}\end{center}\vspace{-2mm}}

% ---------- Boxen ----------
\newenvironment{aufgabe}[2][]{\begin{tcolorbox}[enhanced,breakable,colback=white,colframe=black,boxrule=0.7pt,arc=3mm,left=2mm,right=2mm,top=1mm,bottom=1.2mm,before skip=2mm,after skip=2mm]\textbf{Aufgabe #2}\ifstrempty{#1}{}{ \textit{#1}}\par\vspace{0.3mm}}{\end{tcolorbox}}
\newtcolorbox{merke}[1][Merke]{enhanced,breakable,colback=green!5,colframe=green!40!black,boxrule=0.8pt,arc=2mm,left=2mm,right=2mm,top=1mm,bottom=1.2mm,before skip=2mm,after skip=2mm,title={#1},fonttitle=\bfseries,coltitle=white,toptitle=0.3mm,bottomtitle=0.3mm}
\newtcolorbox{hinweis}{enhanced,breakable,colback=yellow!12,colframe=orange!70!black,boxrule=0.6pt,arc=2mm,left=2mm,right=2mm,top=1mm,bottom=1mm,before skip=2mm,after skip=2mm}

% ---------- Lücken, Linien, Karos ----------
\newcommand{\luecke}[1][5.25cm]{\rule[-0.5ex]{#1}{0.4pt}}
\newcommand{\linien}[1]{\par\vspace{1.5mm}\foreach \i in {1,...,#1}{\noindent\rule[0pt]{\linewidth}{0.3pt}\par\vspace{5.5mm}}\vspace{-3mm}}
\newcommand{\code}[1]{\texttt{#1}}
\newcommand{\eps}{$\varepsilon$}

% Lücke mit Lösung: \lsg[Mindestbreite]{Lösungstext}
%  Die Lösung steht (ohne eigene Höhe/Tiefe) auf derselben Linie wie im Blatt: Das Layout verschiebt sich nicht,
%  solange der Text nicht breiter als die Mindestbreite ist.
\newlength{\lsgbreite}
\newcommand{\lsg}[2][5.25cm]{\leavevmode\ifloesung\settowidth{\lsgbreite}{\,#2\,}\ifdim\lsgbreite<#1\setlength{\lsgbreite}{#1}\fi\rlap{\smash{\raisebox{0.8pt}{\makebox[\lsgbreite][l]{\lsgfarbe\,#2}}}}\luecke[\lsgbreite]\else\luecke[#1]\fi}
% Lücke bis zum Zeilenende (beendet den Absatz): \lsgrest{Lösungstext} bzw. \restzeile
\newcommand{\lsgrest}[1]{\leavevmode\unskip\hspace{1.5mm}\ifloesung\rlap{\smash{\raisebox{0.8pt}{\lsgfarbe\,#1}}}\fi\leaders\hbox{\rule[-0.5ex]{0.5pt}{0.4pt}}\hfill\null\par}
\newcommand{\restzeile}{\lsgrest{}}
% Schreiblinien mit Lösung: \lsglinien{Anzahl}{Lösungstext} (gleiche Höhe wie \linien:
%  je Linie \baselineskip + 5,5 mm + \parskip, ausgemessen)
\newlength{\lsglinienhoehe}
\newcommand{\lsglinien}[2]{\ifloesung\par\vspace{1.5mm}\setlength{\lsglinienhoehe}{\dimexpr #1\dimexpr\baselineskip+5.5mm+\parskip\relax-16.25pt\relax}\noindent\begin{minipage}[t][\lsglinienhoehe][t]{\linewidth}\lsgfarbe #2\end{minipage}\par\vspace{-1.5mm}\else\linien{#1}\fi}
% Beliebiger Inhalt: \lsgoder{nur in der Lösung}{nur im Blatt}
\newcommand{\lsgoder}[2]{\ifloesung\leavevmode#1\else#2\fi}
% Text nur in der Lösung sichtbar, belegt im Blatt aber denselben Platz
\newcommand{\lsgtext}[1]{\leavevmode\ifloesung{\lsgfarbe#1}\else\phantom{#1}\fi}
% Tabellenzelle zum Hineinschreiben (8 mm hoch, ohne Schreiblinie; Lösung mit \lsgtext{…})
\newcommand{\schreibzeile}{\rule[-2.8mm]{0pt}{8mm}}
% Ankreuzfeld: \kreuz{1} = angekreuzt (nur in der Lösung), \kreuz{0} = leer
\newcommand{\kreuz}[1]{\leavevmode\ifloesung\ifnum#1=1 {\lsgfarbe$\boxtimes$}\else$\square$\fi\else$\square$\fi}

% ---------- Syntaxdiagramme (Maße in mm) ----------
%  Terminal rund (orange), Nichtterminal eckig (grün), Linien mit Pfeilen:
%  \draw[sdl,->] (0,0) -- ++(6,0) node[nt,anchor=west] (a) {Augen};   (Pfeil endet am Kasten)
\tikzset{
  sd/.style={x=1mm,y=1mm,>={Stealth[length=1.6mm,width=1.3mm]},font=\small},
  sdl/.style={line width=0.55pt,rounded corners=1.2mm},
  te/.style={draw,rounded corners=2.1mm,minimum height=4.2mm,minimum width=4.2mm,inner sep=1pt,font=\ttfamily\small,fill=orange!18},
  nt/.style={draw,sharp corners,minimum height=4.2mm,inner sep=1.6pt,font=\small,fill=green!12},   % sharp corners: sonst erbt der Kasten die runden Ecken der Linie (sdl)
  sdname/.style={anchor=south west,font=\small\bfseries,inner sep=0pt},
  punkte/.style={font=\small,inner sep=0pt},
  % Ableitungsbaum
  baum/.style={font=\small,edge from parent/.style={draw,line width=0.5pt,-{Stealth[length=1.4mm]}},level distance=8.5mm},
}
% Lösungszeichnung: im Blatt unsichtbar, belegt aber denselben Platz
\ifloesung
  \tikzset{lsgbild/.style={red!75!black,te/.append style={fill=red!5},nt/.append style={fill=red!5}}}
\else
  \tikzset{lsgbild/.style={draw opacity=0,fill opacity=0,text opacity=0}}
\fi
% Karofläche, in der Lösung mit roter Zeichnung: \karolsg{Anzahl Karozeilen}{TikZ-Code}
% Koordinaten der Zeichnung in mm, Ursprung links oben, y nach unten negativ.
\newcommand{\karolsg}[2]{\par\vspace{1mm}\noindent\begin{tikzpicture}
  \draw[step=5mm,gray!45,thin] (0,0) grid (\linewidth-1pt,#1*5mm);
  \useasboundingbox (0,0) rectangle (\linewidth-1pt,#1*5mm);
  \ifloesung\begin{scope}[sd,shift={(0,#1*5mm)},lsgbild]#2\end{scope}\fi
\end{tikzpicture}\par}

% ---------- Bausteine für Syntaxdiagramme (innerhalb tikzpicture[sd]) ----------
%  Jeder Baustein beginnt an einem Punkt (Knoten-Anker oder Koordinate) und
%  legt die Koordinate <name>-o an, an der es weitergeht.
% Aneinanderreihung: \sdnext{name}{von}{te|nt}{Text}   (weiter bei name.east)
\newcommand{\sdnext}[4]{\draw[sdl,->] (#2) -- ++(4,0) node[#3,anchor=west] (#1) {#4};\coordinate (#1-o) at (#1.east);}
% Option: \sdopt{name}{von}{te|nt}{Text}
\newcommand{\sdopt}[4]{\path (#2) ++(4,-5.5) node[#3,anchor=west] (#1) {#4};
  \coordinate (#1-x) at ($(#1.east)+(3.5,0)$); \coordinate (#1-o) at (#1-x |- #2);
  \draw[sdl,->] (#2) |- (#1.west); \draw[sdl,->] (#1.east) -| (#1-o); \draw[sdl] (#2) -- (#1-o);}
% Wiederholung, mindestens einmal: \sdwdh{name}{von}{te|nt}{Text}
\newcommand{\sdwdh}[4]{\draw[sdl,->] (#2) -- ++(6,0) node[#3,anchor=west] (#1) {#4};
  \coordinate (#1-o) at ($(#1.east)+(4,0)$);
  \draw[sdl] (#1.east) -- (#1-o); \draw[sdl,->] (#1-o) -- ++(0,-5.5) -| ($(#2)+(2.5,0)$);}
% Wiederholung, beliebig oft (auch keinmal): \sdwdhn{name}{von}{te|nt}{Text}
\newcommand{\sdwdhn}[4]{\path (#2) ++(6,-5.5) node[#3,anchor=west] (#1) {#4};
  \coordinate (#1-x) at ($(#1.east)+(4,0)$); \coordinate (#1-o) at (#1-x |- #2);
  \draw[sdl] (#2) -- (#1-o); \draw[sdl,->] (#1-o) |- (#1.east); \draw[sdl,->] (#1.west) -| ($(#2)+(2.5,0)$);}
% Alternative: \sdalt{name}{von}{Breite bis zum Zusammenführen}{Zeilenabstand}{te/a, nt/Ziffer, pk/, ...}
%  pk/ setzt senkrechte Pünktchen (…) statt eines Zweigs
\newcommand{\sdalt}[5]{\coordinate (#1-o) at ($(#2)+(#3,0)$);
  \foreach \typ/\txt [count=\i from 0] in {#5}{%
    \ifdefstring{\typ}{pk}{\node[punkte] at ($(#2)+(6.1,-\i*#4+0.4)$) {$\vdots$};}{%
      \ifdefstring{\typ}{te}{\tikzset{sdakt/.style={te}}}{\tikzset{sdakt/.style={nt}}}%
      \node[sdakt,anchor=west] (#1\i) at ($(#2)+(4,-\i*#4)$) {\txt};
      \ifnum\i=0\relax \draw[sdl,->] (#2) -- (#1\i.west); \draw[sdl] (#1\i.east) -- (#1-o);
      \else \draw[sdl,->] (#2) |- (#1\i.west); \draw[sdl,->] (#1\i.east) -| (#1-o);\fi}}}
% Ende eines Diagramms
\newcommand{\sdende}[1]{\draw[sdl,->] (#1) -- ++(5,0);}

% ---------- Die Smiley-Grammatik (Blatt 1.3 und 1.4) ----------
\newcommand{\sdSmileyfolge}{\begin{tikzpicture}[sd,baseline=0mm]
  \node[sdname] at (0,3.3) {Smileyfolge}; \coordinate (s) at (0,0);
  \sdwdh{sm}{s}{nt}{Smiley}\sdende{sm-o}\end{tikzpicture}}
\newcommand{\sdSmiley}{\begin{tikzpicture}[sd,baseline=0mm]
  \node[sdname] at (0,3.3) {Smiley}; \coordinate (s) at (0,0);
  \sdnext{au}{s}{nt}{Augen}\sdnext{mi}{au-o}{nt}{Mitte}\sdnext{mu}{mi-o}{nt}{Mund}\sdende{mu-o}\end{tikzpicture}}
\newcommand{\sdAugen}{\begin{tikzpicture}[sd,baseline=0mm]
  \node[sdname] at (0,3.3) {Augen}; \coordinate (s) at (0,0);
  \sdalt{a}{s}{12}{5.5}{te/:, te/;}\sdende{a-o}\end{tikzpicture}}
\newcommand{\sdMitte}{\begin{tikzpicture}[sd,baseline=0mm]
  \node[sdname] at (0,3.3) {Mitte}; \coordinate (s) at (0,0);
  \sdopt{n}{s}{te}{-}\sdende{n-o}\end{tikzpicture}}
\newcommand{\sdMund}{\begin{tikzpicture}[sd,baseline=0mm]
  \node[sdname] at (0,3.3) {Mund}; \coordinate (s) at (0,0);
  \sdalt{m}{s}{12}{5.5}{te/), te/{(}, te/D, te/P, te/*}\sdende{m-o}\end{tikzpicture}}

% ---------- Ergänzungen für Programmier-Blätter (ab 1.6; aus informatik/10/java/arbeitsblaetter/ab-vorlage.tex) ----------
%  Nur angehängt: Die Makros oben bleiben unverändert, damit die PDFs von 1.1 bis 1.5 gleich bleiben.
%  \lueckentext   am Anfang eines Absatzes/Kastens mit Lücken: Zeilenabstand 1,5, Flattersatz
%  \nr{1}         Schrittnummer ① wie auf der Website und im Code
%  java           Listing:  \begin{java}[numbers=left] … \end{java}
%  javaluecke     Lückencode mit größerem Zeilenabstand, Lücken als (*\lsg[6.75cm]{Lösung}*)
\newcommand{\lueckentext}{\linespread{1.5}\selectfont\raggedright}
\newcommand{\nr}[1]{\ding{\numexpr171+#1\relax}}
\usepackage{listings}
\lstset{language=Java,basicstyle=\ttfamily\small,keywordstyle=\bfseries,commentstyle=\color{gray!80!black},stringstyle=\color{black},showstringspaces=false,columns=flexible,keepspaces=true,tabsize=3,frame=single,framerule=0.4pt,rulecolor=\color{black},xleftmargin=2mm,framexleftmargin=1mm,aboveskip=1.5mm,belowskip=1.5mm,breaklines=true,numberstyle=\tiny\color{gray},numbersep=3mm,escapeinside={(*}{*)},
 literate={ä}{{\"a}}1 {ö}{{\"o}}1 {ü}{{\"u}}1 {Ä}{{\"A}}1 {Ö}{{\"O}}1 {Ü}{{\"U}}1 {ß}{{\ss}}1 {–}{{--}}1 {…}{{\dots}}1}
\lstnewenvironment{java}[1][]{\lstset{#1}}{}
\lstnewenvironment{javaluecke}[1][]{\lstset{basicstyle=\linespread{1.5}\small\ttfamily,#1}}{}

%  Übersetzen mit XeLaTeX, LuaLaTeX, pdfLaTeX oder Tectonic.
%
%  Blatt:   tectonic ab-1-2-formale-sprachen.tex
%  Lösung:  tectonic ab-1-2-formale-sprachen-lsg.tex
%           (Datei enthält nur: \def\mitloesung{}\documentclass[11pt,a4paper]{article}
% =====================================================================
%  Gemeinsame Vorlage der Arbeitsblätter – 1. Formale Sprachen (Q13)
%  (gleiches Layout wie informatik/13/material/grundwissen/ab-vorlage.tex
%  und informatik/12/rekursive-datenstrukturen/arbeitsblaetter/ab-vorlage.tex)
%  Einbinden in jedem Blatt direkt nach \documentclass: \input{ab-vorlage}
%  Übersetzen mit XeLaTeX, LuaLaTeX, pdfLaTeX oder Tectonic.
%
%  Blatt:   tectonic ab-1-2-formale-sprachen.tex
%  Lösung:  tectonic ab-1-2-formale-sprachen-lsg.tex
%           (Datei enthält nur: \def\mitloesung{}\input{ab-1-2-formale-sprachen};
%            füllt alle Lücken rot aus; siehe \lsg, \lsglinien, \lsgoder, lsgbild, \karolsg)
% =====================================================================
\usepackage[a4paper,left=16mm,right=16mm,top=23mm,bottom=13mm,headheight=12mm,headsep=3mm,footskip=6mm]{geometry}
\usepackage{iftex}
\ifPDFTeX
  \usepackage[T1]{fontenc}
  \usepackage[utf8]{inputenc}
  \usepackage{helvet}
  \usepackage[scaled=0.86]{beramono}
  \renewcommand{\familydefault}{\sfdefault}
\else
  \usepackage{fontspec}
  \setmainfont{texgyreheros}[Extension=.otf,UprightFont=*-regular,BoldFont=*-bold,ItalicFont=*-italic,BoldItalicFont=*-bolditalic]
  \setmonofont{DejaVuSansMono}[Extension=.ttf,UprightFont=*,BoldFont=*-Bold,ItalicFont=*-Oblique,BoldItalicFont=*-BoldOblique,Scale=0.86]
\fi
\usepackage[ngerman,shorthands=off]{babel}
\usepackage{amssymb}
\usepackage{xcolor,graphicx,tabularx,array,enumitem,fancyhdr,tikz,multicol,calc,etoolbox,pifont}
\usepackage[most]{tcolorbox}
\usetikzlibrary{arrows.meta,positioning,calc}
\usepackage[hidelinks]{hyperref}
\setlength{\parindent}{0pt}
\setlength{\parskip}{3pt}
\setlist{nosep,leftmargin=*}

% ---------- Lösungsschalter ----------
\newif\ifloesung
\ifdefined\mitloesung\loesungtrue\fi
\newcommand{\lsgfarbe}{\color{red!75!black}}

% ---------- Kopf- und Fußzeile (einheitliche Form) ----------
\newcommand{\lehrkraft}{Hau}
\newcommand{\themenbereich}{1. Formale Sprachen}
\newcommand{\abnummer}{}
\pagestyle{fancy}\fancyhf{}
\renewcommand{\headrulewidth}{0.4pt}
\fancyhead[L]{\small Klasse 13\\Datum:}
\fancyhead[C]{\small Themenbereich:\\\themenbereich}
\fancyhead[R]{\small \lehrkraft\ifloesung\\{\lsgfarbe\bfseries Lösung}\fi}
\fancyfoot[C]{\scriptsize Informatik 13 gA\,·\,Blatt \abnummer\ifloesung\ (Lösung)\fi\,·\,Seite \thepage}
\newcommand{\abtitel}[2]{\renewcommand{\abnummer}{#1}\begin{center}\Large\bfseries #1\quad\underline{#2}\end{center}\vspace{-2mm}}

% ---------- Boxen ----------
\newenvironment{aufgabe}[2][]{\begin{tcolorbox}[enhanced,breakable,colback=white,colframe=black,boxrule=0.7pt,arc=3mm,left=2mm,right=2mm,top=1mm,bottom=1.2mm,before skip=2mm,after skip=2mm]\textbf{Aufgabe #2}\ifstrempty{#1}{}{ \textit{#1}}\par\vspace{0.3mm}}{\end{tcolorbox}}
\newtcolorbox{merke}[1][Merke]{enhanced,breakable,colback=green!5,colframe=green!40!black,boxrule=0.8pt,arc=2mm,left=2mm,right=2mm,top=1mm,bottom=1.2mm,before skip=2mm,after skip=2mm,title={#1},fonttitle=\bfseries,coltitle=white,toptitle=0.3mm,bottomtitle=0.3mm}
\newtcolorbox{hinweis}{enhanced,breakable,colback=yellow!12,colframe=orange!70!black,boxrule=0.6pt,arc=2mm,left=2mm,right=2mm,top=1mm,bottom=1mm,before skip=2mm,after skip=2mm}

% ---------- Lücken, Linien, Karos ----------
\newcommand{\luecke}[1][5.25cm]{\rule[-0.5ex]{#1}{0.4pt}}
\newcommand{\linien}[1]{\par\vspace{1.5mm}\foreach \i in {1,...,#1}{\noindent\rule[0pt]{\linewidth}{0.3pt}\par\vspace{5.5mm}}\vspace{-3mm}}
\newcommand{\code}[1]{\texttt{#1}}
\newcommand{\eps}{$\varepsilon$}

% Lücke mit Lösung: \lsg[Mindestbreite]{Lösungstext}
%  Die Lösung steht (ohne eigene Höhe/Tiefe) auf derselben Linie wie im Blatt: Das Layout verschiebt sich nicht,
%  solange der Text nicht breiter als die Mindestbreite ist.
\newlength{\lsgbreite}
\newcommand{\lsg}[2][5.25cm]{\leavevmode\ifloesung\settowidth{\lsgbreite}{\,#2\,}\ifdim\lsgbreite<#1\setlength{\lsgbreite}{#1}\fi\rlap{\smash{\raisebox{0.8pt}{\makebox[\lsgbreite][l]{\lsgfarbe\,#2}}}}\luecke[\lsgbreite]\else\luecke[#1]\fi}
% Lücke bis zum Zeilenende (beendet den Absatz): \lsgrest{Lösungstext} bzw. \restzeile
\newcommand{\lsgrest}[1]{\leavevmode\unskip\hspace{1.5mm}\ifloesung\rlap{\smash{\raisebox{0.8pt}{\lsgfarbe\,#1}}}\fi\leaders\hbox{\rule[-0.5ex]{0.5pt}{0.4pt}}\hfill\null\par}
\newcommand{\restzeile}{\lsgrest{}}
% Schreiblinien mit Lösung: \lsglinien{Anzahl}{Lösungstext} (gleiche Höhe wie \linien:
%  je Linie \baselineskip + 5,5 mm + \parskip, ausgemessen)
\newlength{\lsglinienhoehe}
\newcommand{\lsglinien}[2]{\ifloesung\par\vspace{1.5mm}\setlength{\lsglinienhoehe}{\dimexpr #1\dimexpr\baselineskip+5.5mm+\parskip\relax-16.25pt\relax}\noindent\begin{minipage}[t][\lsglinienhoehe][t]{\linewidth}\lsgfarbe #2\end{minipage}\par\vspace{-1.5mm}\else\linien{#1}\fi}
% Beliebiger Inhalt: \lsgoder{nur in der Lösung}{nur im Blatt}
\newcommand{\lsgoder}[2]{\ifloesung\leavevmode#1\else#2\fi}
% Text nur in der Lösung sichtbar, belegt im Blatt aber denselben Platz
\newcommand{\lsgtext}[1]{\leavevmode\ifloesung{\lsgfarbe#1}\else\phantom{#1}\fi}
% Tabellenzelle zum Hineinschreiben (8 mm hoch, ohne Schreiblinie; Lösung mit \lsgtext{…})
\newcommand{\schreibzeile}{\rule[-2.8mm]{0pt}{8mm}}
% Ankreuzfeld: \kreuz{1} = angekreuzt (nur in der Lösung), \kreuz{0} = leer
\newcommand{\kreuz}[1]{\leavevmode\ifloesung\ifnum#1=1 {\lsgfarbe$\boxtimes$}\else$\square$\fi\else$\square$\fi}

% ---------- Syntaxdiagramme (Maße in mm) ----------
%  Terminal rund (orange), Nichtterminal eckig (grün), Linien mit Pfeilen:
%  \draw[sdl,->] (0,0) -- ++(6,0) node[nt,anchor=west] (a) {Augen};   (Pfeil endet am Kasten)
\tikzset{
  sd/.style={x=1mm,y=1mm,>={Stealth[length=1.6mm,width=1.3mm]},font=\small},
  sdl/.style={line width=0.55pt,rounded corners=1.2mm},
  te/.style={draw,rounded corners=2.1mm,minimum height=4.2mm,minimum width=4.2mm,inner sep=1pt,font=\ttfamily\small,fill=orange!18},
  nt/.style={draw,sharp corners,minimum height=4.2mm,inner sep=1.6pt,font=\small,fill=green!12},   % sharp corners: sonst erbt der Kasten die runden Ecken der Linie (sdl)
  sdname/.style={anchor=south west,font=\small\bfseries,inner sep=0pt},
  punkte/.style={font=\small,inner sep=0pt},
  % Ableitungsbaum
  baum/.style={font=\small,edge from parent/.style={draw,line width=0.5pt,-{Stealth[length=1.4mm]}},level distance=8.5mm},
}
% Lösungszeichnung: im Blatt unsichtbar, belegt aber denselben Platz
\ifloesung
  \tikzset{lsgbild/.style={red!75!black,te/.append style={fill=red!5},nt/.append style={fill=red!5}}}
\else
  \tikzset{lsgbild/.style={draw opacity=0,fill opacity=0,text opacity=0}}
\fi
% Karofläche, in der Lösung mit roter Zeichnung: \karolsg{Anzahl Karozeilen}{TikZ-Code}
% Koordinaten der Zeichnung in mm, Ursprung links oben, y nach unten negativ.
\newcommand{\karolsg}[2]{\par\vspace{1mm}\noindent\begin{tikzpicture}
  \draw[step=5mm,gray!45,thin] (0,0) grid (\linewidth-1pt,#1*5mm);
  \useasboundingbox (0,0) rectangle (\linewidth-1pt,#1*5mm);
  \ifloesung\begin{scope}[sd,shift={(0,#1*5mm)},lsgbild]#2\end{scope}\fi
\end{tikzpicture}\par}

% ---------- Bausteine für Syntaxdiagramme (innerhalb tikzpicture[sd]) ----------
%  Jeder Baustein beginnt an einem Punkt (Knoten-Anker oder Koordinate) und
%  legt die Koordinate <name>-o an, an der es weitergeht.
% Aneinanderreihung: \sdnext{name}{von}{te|nt}{Text}   (weiter bei name.east)
\newcommand{\sdnext}[4]{\draw[sdl,->] (#2) -- ++(4,0) node[#3,anchor=west] (#1) {#4};\coordinate (#1-o) at (#1.east);}
% Option: \sdopt{name}{von}{te|nt}{Text}
\newcommand{\sdopt}[4]{\path (#2) ++(4,-5.5) node[#3,anchor=west] (#1) {#4};
  \coordinate (#1-x) at ($(#1.east)+(3.5,0)$); \coordinate (#1-o) at (#1-x |- #2);
  \draw[sdl,->] (#2) |- (#1.west); \draw[sdl,->] (#1.east) -| (#1-o); \draw[sdl] (#2) -- (#1-o);}
% Wiederholung, mindestens einmal: \sdwdh{name}{von}{te|nt}{Text}
\newcommand{\sdwdh}[4]{\draw[sdl,->] (#2) -- ++(6,0) node[#3,anchor=west] (#1) {#4};
  \coordinate (#1-o) at ($(#1.east)+(4,0)$);
  \draw[sdl] (#1.east) -- (#1-o); \draw[sdl,->] (#1-o) -- ++(0,-5.5) -| ($(#2)+(2.5,0)$);}
% Wiederholung, beliebig oft (auch keinmal): \sdwdhn{name}{von}{te|nt}{Text}
\newcommand{\sdwdhn}[4]{\path (#2) ++(6,-5.5) node[#3,anchor=west] (#1) {#4};
  \coordinate (#1-x) at ($(#1.east)+(4,0)$); \coordinate (#1-o) at (#1-x |- #2);
  \draw[sdl] (#2) -- (#1-o); \draw[sdl,->] (#1-o) |- (#1.east); \draw[sdl,->] (#1.west) -| ($(#2)+(2.5,0)$);}
% Alternative: \sdalt{name}{von}{Breite bis zum Zusammenführen}{Zeilenabstand}{te/a, nt/Ziffer, pk/, ...}
%  pk/ setzt senkrechte Pünktchen (…) statt eines Zweigs
\newcommand{\sdalt}[5]{\coordinate (#1-o) at ($(#2)+(#3,0)$);
  \foreach \typ/\txt [count=\i from 0] in {#5}{%
    \ifdefstring{\typ}{pk}{\node[punkte] at ($(#2)+(6.1,-\i*#4+0.4)$) {$\vdots$};}{%
      \ifdefstring{\typ}{te}{\tikzset{sdakt/.style={te}}}{\tikzset{sdakt/.style={nt}}}%
      \node[sdakt,anchor=west] (#1\i) at ($(#2)+(4,-\i*#4)$) {\txt};
      \ifnum\i=0\relax \draw[sdl,->] (#2) -- (#1\i.west); \draw[sdl] (#1\i.east) -- (#1-o);
      \else \draw[sdl,->] (#2) |- (#1\i.west); \draw[sdl,->] (#1\i.east) -| (#1-o);\fi}}}
% Ende eines Diagramms
\newcommand{\sdende}[1]{\draw[sdl,->] (#1) -- ++(5,0);}

% ---------- Die Smiley-Grammatik (Blatt 1.3 und 1.4) ----------
\newcommand{\sdSmileyfolge}{\begin{tikzpicture}[sd,baseline=0mm]
  \node[sdname] at (0,3.3) {Smileyfolge}; \coordinate (s) at (0,0);
  \sdwdh{sm}{s}{nt}{Smiley}\sdende{sm-o}\end{tikzpicture}}
\newcommand{\sdSmiley}{\begin{tikzpicture}[sd,baseline=0mm]
  \node[sdname] at (0,3.3) {Smiley}; \coordinate (s) at (0,0);
  \sdnext{au}{s}{nt}{Augen}\sdnext{mi}{au-o}{nt}{Mitte}\sdnext{mu}{mi-o}{nt}{Mund}\sdende{mu-o}\end{tikzpicture}}
\newcommand{\sdAugen}{\begin{tikzpicture}[sd,baseline=0mm]
  \node[sdname] at (0,3.3) {Augen}; \coordinate (s) at (0,0);
  \sdalt{a}{s}{12}{5.5}{te/:, te/;}\sdende{a-o}\end{tikzpicture}}
\newcommand{\sdMitte}{\begin{tikzpicture}[sd,baseline=0mm]
  \node[sdname] at (0,3.3) {Mitte}; \coordinate (s) at (0,0);
  \sdopt{n}{s}{te}{-}\sdende{n-o}\end{tikzpicture}}
\newcommand{\sdMund}{\begin{tikzpicture}[sd,baseline=0mm]
  \node[sdname] at (0,3.3) {Mund}; \coordinate (s) at (0,0);
  \sdalt{m}{s}{12}{5.5}{te/), te/{(}, te/D, te/P, te/*}\sdende{m-o}\end{tikzpicture}}

% ---------- Ergänzungen für Programmier-Blätter (ab 1.6; aus informatik/10/java/arbeitsblaetter/ab-vorlage.tex) ----------
%  Nur angehängt: Die Makros oben bleiben unverändert, damit die PDFs von 1.1 bis 1.5 gleich bleiben.
%  \lueckentext   am Anfang eines Absatzes/Kastens mit Lücken: Zeilenabstand 1,5, Flattersatz
%  \nr{1}         Schrittnummer ① wie auf der Website und im Code
%  java           Listing:  \begin{java}[numbers=left] … \end{java}
%  javaluecke     Lückencode mit größerem Zeilenabstand, Lücken als (*\lsg[6.75cm]{Lösung}*)
\newcommand{\lueckentext}{\linespread{1.5}\selectfont\raggedright}
\newcommand{\nr}[1]{\ding{\numexpr171+#1\relax}}
\usepackage{listings}
\lstset{language=Java,basicstyle=\ttfamily\small,keywordstyle=\bfseries,commentstyle=\color{gray!80!black},stringstyle=\color{black},showstringspaces=false,columns=flexible,keepspaces=true,tabsize=3,frame=single,framerule=0.4pt,rulecolor=\color{black},xleftmargin=2mm,framexleftmargin=1mm,aboveskip=1.5mm,belowskip=1.5mm,breaklines=true,numberstyle=\tiny\color{gray},numbersep=3mm,escapeinside={(*}{*)},
 literate={ä}{{\"a}}1 {ö}{{\"o}}1 {ü}{{\"u}}1 {Ä}{{\"A}}1 {Ö}{{\"O}}1 {Ü}{{\"U}}1 {ß}{{\ss}}1 {–}{{--}}1 {…}{{\dots}}1}
\lstnewenvironment{java}[1][]{\lstset{#1}}{}
\lstnewenvironment{javaluecke}[1][]{\lstset{basicstyle=\linespread{1.5}\small\ttfamily,#1}}{}

% Blatt:   tectonic ab-1-2-formale-sprachen.tex
% Lösung:  tectonic ab-1-2-formale-sprachen-lsg.tex
% Übungsfassung: Die Begriffe (Alphabet, Wort, Länge, ε, A*, L ⊆ A*) führt Blatt 1.1 ein, hier wird nur geübt.
% Progression: Begriffe sicher anwenden (ε!) -> Wörter zählen |A|^n, endliche Sprache in unendlichem A*
%              -> Sprachen durch Eigenschaften festlegen -> Alltag (Syntax, nicht Semantik)
%              -> Smileys in Worten beschreiben: mehrdeutig -> Regeln nötig (Brücke zu 1.3)
\newcommand{\zeile}{\rule[-2mm]{0pt}{6.6mm}}
\newcommand{\hoch}{\rule[-1.6mm]{0pt}{6mm}}
\begin{document}
\abtitel{1.2}{Formale Sprachen -- Übung}

Die Begriffe kennst du aus 1.1 (Merkkasten \emph{Formale Sprachen}). Auf diesem Blatt übst du sie mit drei Alphabeten:\\[1mm]
\centerline{$A_1 = \{\code{0}, \code{1}\}$\qquad $A_2 = \{\code{a}, \code{b}, \code{c}, \dots, \code{z}\}$\qquad $A_3 = \{\code{:}, \code{;}, \code{-}, \code{)}, \code{(}, \code{D}, \code{P}, \code{*}\}$}

\begin{aufgabe}[Kurz geprüft]{1}
\textbf{a)} Kreuze alle richtigen Aussagen an.\\[1mm]
\renewcommand{\arraystretch}{1.25}
\begin{tabularx}{\linewidth}{@{}X@{\quad}X@{}}
\kreuz{1} $\code{0} \in A_1$ & \kreuz{0} $\varepsilon \in A_1$\\
\kreuz{0} \code{ab} ist ein Wort über $A_1$ & \kreuz{1} $\varepsilon \in A_1^*$\\
\kreuz{0} $A_1^*$ ist endlich, weil $A_1$ endlich ist. & \kreuz{1} \code{:-)} ist ein Wort der Länge 3 über $A_3$.\\
\kreuz{1} Jedes Zeichen aus $A$ ist auch ein Wort über $A$. & \kreuz{0} Das leere Wort $\varepsilon$ ist ein Leerzeichen.\\
\end{tabularx}\\[1mm]
\textbf{b)} Es gibt genau ein Wort, das über \emph{jedem} Alphabet gebildet werden kann. Welches? Warum?\\[0.6mm]
\lsg[\linewidth]{Das leere Wort \eps: Es enthält kein Zeichen, also auch keines, das im Alphabet fehlen könnte.}
\end{aufgabe}

\begin{aufgabe}[Wie viele Wörter gibt es?]{2}
\textbf{a)} Schreibe alle Wörter über $A_1$ der Länge 0 bis 3 auf und zähle. Zähle auch für $A_3$ und allgemein.\\[1mm]
\renewcommand{\arraystretch}{1.1}
\begin{tabularx}{\linewidth}{|p{4.4cm}|p{0.8cm}|p{1.3cm}|p{2.6cm}|X|p{1.1cm}|}\hline
\textbf{Länge $n$} & \textbf{0} & \textbf{1} & \textbf{2} & \textbf{3} & \textbf{$n$}\\\hline
Wörter über $A_1$\rule[-6.5mm]{0pt}{11mm} & \lsgoder{\lsgfarbe \eps}{} & \lsgoder{\lsgfarbe \code{0}, \code{1}}{} & \lsgoder{\lsgfarbe \code{00}, \code{01}, \code{10}, \code{11}}{} & \lsgoder{\lsgfarbe \code{000}, \code{001}, \code{010}, \code{011}, \code{100}, \code{101}, \code{110}, \code{111}}{} & \\\hline
Anzahl über $A_1$\hoch & \lsgoder{\lsgfarbe 1}{} & \lsgoder{\lsgfarbe 2}{} & \lsgoder{\lsgfarbe 4}{} & \lsgoder{\lsgfarbe 8}{} & \lsgoder{\lsgfarbe $2^n$}{}\\\hline
Anzahl über $A_3$\hoch & \lsgoder{\lsgfarbe 1}{} & \lsgoder{\lsgfarbe 8}{} & \lsgoder{\lsgfarbe 64}{} & \lsgoder{\lsgfarbe 512}{} & \lsgoder{\lsgfarbe $8^n$}{}\\\hline
Anzahl bei $k$ Zeichen in $A$\hoch & \lsgoder{\lsgfarbe 1}{} & \lsgoder{\lsgfarbe $k$}{} & \lsgoder{\lsgfarbe $k^2$}{} & \lsgoder{\lsgfarbe $k^3$}{} & \lsgoder{\lsgfarbe $k^n$}{}\\\hline
\end{tabularx}\\[1.2mm]
\textbf{b)} Wie viele Farbcodes wie \code{\#9900BB} gibt es (1.1, Aufgabe 6)? Ist diese Sprache endlich oder unendlich?\\[0.6mm]
\lsg[\linewidth]{$16^6 = 16\,777\,216$ Farbcodes: Die Sprache ist endlich, obwohl $A_\text{Hex}^*$ unendlich ist.}
\end{aufgabe}

\begin{aufgabe}[Sprachen über $A_2$]{3}
\begin{minipage}[t]{0.585\linewidth}
\vspace{0pt}
$L_1$ = alle Wörter aus $A_2^*$ mit Länge kleiner gleich 4\\
$L_2$ = alle Wörter aus $A_2^*$, die mit \code{b} beginnen\\
$L_3$ = alle Wörter aus $A_2^*$, die das Teilwort \code{inf} enthalten\\[1mm]
\textbf{a)} Kreuze an, zu welcher Sprache das Wort gehört.\\[1mm]
\textbf{b)} Welche der drei Sprachen sind endlich? \lsg[2.7cm]{nur $L_1$}\\[1mm]
\textbf{c)} Finde ein Wort, das zu $L_1$, $L_2$ \emph{und} $L_3$ gehört:\\[0.6mm] \lsg[4.8cm]{\code{binf}}\\[1.2mm]
Gibt es mehrere? \lsg[3.9cm]{nein, nur \code{binf}}
\end{minipage}\hfill
\begin{minipage}[t]{0.39\linewidth}
\vspace{0pt}
\renewcommand{\arraystretch}{1.12}
\begin{tabularx}{\linewidth}{|X|c|c|c|}\hline
\textbf{Wort} & $L_1$ & $L_2$ & $L_3$\\\hline
\code{baum} & \kreuz{1} & \kreuz{1} & \kreuz{0}\\\hline
\code{informatik} & \kreuz{0} & \kreuz{0} & \kreuz{1}\\\hline
\eps & \kreuz{1} & \kreuz{0} & \kreuz{0}\\\hline
\code{inf} & \kreuz{1} & \kreuz{0} & \kreuz{1}\\\hline
\code{binfo} & \kreuz{0} & \kreuz{1} & \kreuz{1}\\\hline
\end{tabularx}
\end{minipage}
\end{aufgabe}

\newpage
\begin{aufgabe}[Formale Sprachen im Alltag]{4}
\textbf{a)} Ergänze die Tabelle.\\[1mm]
\renewcommand{\arraystretch}{1.1}
\begin{tabularx}{\linewidth}{|p{5.3cm}|X|p{2.4cm}|p{4.5cm}|}\hline
\textbf{Sprache} & \textbf{Alphabet} & \textbf{Wort aus $L$} & \textbf{Wort aus $A^*$, nicht aus $L$}\\\hline
Postleitzahlen\zeile & \lsgoder{\lsgfarbe \{0, 1, \dots, 9\}}{} & \code{83278} & \lsgoder{\lsgfarbe \code{824377} (6 Ziffern)}{}\\\hline
Binärzahlen ohne führende Null\zeile & \lsgoder{\lsgfarbe \{0, 1\}}{} & \lsgoder{\lsgfarbe \code{1011}}{} & \lsgoder{\lsgfarbe \code{0011}}{}\\\hline
Datum (TT.MM.JJJJ)\zeile & \lsgoder{\lsgfarbe \{0, 1, \dots, 9, .\}}{} & \lsgoder{\lsgfarbe \code{29.09.2026}}{} & \lsgoder{\lsgfarbe \code{29..092026}}{}\\\hline
\end{tabularx}\\[1.2mm]
\textbf{b)} Die Postleitzahl \code{99999} ist nicht vergeben. Gehört sie trotzdem zur Sprache „alle Wörter aus fünf Ziffern“? Was folgt daraus für Syntax und Semantik?
\lsglinien{2}{Ja, \code{99999} ist syntaktisch korrekt. Eine formale Sprache legt nur die \textbf{Syntax} fest; ob das Wort auch etwas bedeutet (\textbf{Semantik}), prüft sie nicht.}
\end{aufgabe}

\begin{aufgabe}[Die Sprache der Smileys]{5}
Im Chat werden Smileys aus den Zeichen von $A_3$ gebildet. Sie bilden eine Sprache $L_\text{Smiley} \subseteq A_3^*$.\\[1mm]
\textbf{a)} Kreuze an, welche Wörter du als Smiley erkennst.\\[1mm]
\kreuz{1} \code{:-)}\qquad \kreuz{1} \code{;D}\qquad \kreuz{1} \code{:-P}\qquad \kreuz{0} \code{)-:}\qquad \kreuz{0} \code{:-}\qquad \kreuz{1} \code{;*}\qquad \kreuz{0} \code{:--)}\qquad \kreuz{0} \code{DP}\qquad \kreuz{1} \code{:(}\\[1mm]
\textbf{b)} Beschreibe in Worten möglichst genau, welche Wörter aus $A_3^*$ zu $L_\text{Smiley}$ gehören.
\lsglinien{3}{Ein Smiley beginnt mit Augen (\code{:} oder \code{;}). Danach kann eine Nase \code{-} folgen, höchstens eine. Am Ende steht genau ein Mund: \code{)}, \code{(}, \code{D}, \code{P} oder \code{*}. Sonst folgt nichts.}
\textbf{c)} Vergleiche mit deinem Nachbarn: Gibt es Wörter, bei denen eure Beschreibungen verschieden entscheiden oder gar nicht entscheiden? Notiere ein Beispiel und woran es liegt.
\lsglinien{2}{z.\,B. \code{)-:} (gespiegelt?), \code{:--)} (wie viele Nasen?), \code{:-):-)} (mehrere Smileys?). Umgangssprache ist ungenau und mehrdeutig; man übersieht leicht Fälle.}
\end{aufgabe}

\begin{merke}[Erkenntnis]
Eine Beschreibung in Umgangssprache legt eine Sprache oft \lsg[4.2cm]{nicht eindeutig} fest. Nötig sind genaue \textbf{Regeln}, nach denen man für jedes Wort entscheiden kann, ob es zur Sprache gehört.
\end{merke}
\end{document}
;
%            füllt alle Lücken rot aus; siehe \lsg, \lsglinien, \lsgoder, lsgbild, \karolsg)
% =====================================================================
\usepackage[a4paper,left=16mm,right=16mm,top=23mm,bottom=13mm,headheight=12mm,headsep=3mm,footskip=6mm]{geometry}
\usepackage{iftex}
\ifPDFTeX
  \usepackage[T1]{fontenc}
  \usepackage[utf8]{inputenc}
  \usepackage{helvet}
  \usepackage[scaled=0.86]{beramono}
  \renewcommand{\familydefault}{\sfdefault}
\else
  \usepackage{fontspec}
  \setmainfont{texgyreheros}[Extension=.otf,UprightFont=*-regular,BoldFont=*-bold,ItalicFont=*-italic,BoldItalicFont=*-bolditalic]
  \setmonofont{DejaVuSansMono}[Extension=.ttf,UprightFont=*,BoldFont=*-Bold,ItalicFont=*-Oblique,BoldItalicFont=*-BoldOblique,Scale=0.86]
\fi
\usepackage[ngerman,shorthands=off]{babel}
\usepackage{amssymb}
\usepackage{xcolor,graphicx,tabularx,array,enumitem,fancyhdr,tikz,multicol,calc,etoolbox,pifont}
\usepackage[most]{tcolorbox}
\usetikzlibrary{arrows.meta,positioning,calc}
\usepackage[hidelinks]{hyperref}
\setlength{\parindent}{0pt}
\setlength{\parskip}{3pt}
\setlist{nosep,leftmargin=*}

% ---------- Lösungsschalter ----------
\newif\ifloesung
\ifdefined\mitloesung\loesungtrue\fi
\newcommand{\lsgfarbe}{\color{red!75!black}}

% ---------- Kopf- und Fußzeile (einheitliche Form) ----------
\newcommand{\lehrkraft}{Hau}
\newcommand{\themenbereich}{1. Formale Sprachen}
\newcommand{\abnummer}{}
\pagestyle{fancy}\fancyhf{}
\renewcommand{\headrulewidth}{0.4pt}
\fancyhead[L]{\small Klasse 13\\Datum:}
\fancyhead[C]{\small Themenbereich:\\\themenbereich}
\fancyhead[R]{\small \lehrkraft\ifloesung\\{\lsgfarbe\bfseries Lösung}\fi}
\fancyfoot[C]{\scriptsize Informatik 13 gA\,·\,Blatt \abnummer\ifloesung\ (Lösung)\fi\,·\,Seite \thepage}
\newcommand{\abtitel}[2]{\renewcommand{\abnummer}{#1}\begin{center}\Large\bfseries #1\quad\underline{#2}\end{center}\vspace{-2mm}}

% ---------- Boxen ----------
\newenvironment{aufgabe}[2][]{\begin{tcolorbox}[enhanced,breakable,colback=white,colframe=black,boxrule=0.7pt,arc=3mm,left=2mm,right=2mm,top=1mm,bottom=1.2mm,before skip=2mm,after skip=2mm]\textbf{Aufgabe #2}\ifstrempty{#1}{}{ \textit{#1}}\par\vspace{0.3mm}}{\end{tcolorbox}}
\newtcolorbox{merke}[1][Merke]{enhanced,breakable,colback=green!5,colframe=green!40!black,boxrule=0.8pt,arc=2mm,left=2mm,right=2mm,top=1mm,bottom=1.2mm,before skip=2mm,after skip=2mm,title={#1},fonttitle=\bfseries,coltitle=white,toptitle=0.3mm,bottomtitle=0.3mm}
\newtcolorbox{hinweis}{enhanced,breakable,colback=yellow!12,colframe=orange!70!black,boxrule=0.6pt,arc=2mm,left=2mm,right=2mm,top=1mm,bottom=1mm,before skip=2mm,after skip=2mm}

% ---------- Lücken, Linien, Karos ----------
\newcommand{\luecke}[1][5.25cm]{\rule[-0.5ex]{#1}{0.4pt}}
\newcommand{\linien}[1]{\par\vspace{1.5mm}\foreach \i in {1,...,#1}{\noindent\rule[0pt]{\linewidth}{0.3pt}\par\vspace{5.5mm}}\vspace{-3mm}}
\newcommand{\code}[1]{\texttt{#1}}
\newcommand{\eps}{$\varepsilon$}

% Lücke mit Lösung: \lsg[Mindestbreite]{Lösungstext}
%  Die Lösung steht (ohne eigene Höhe/Tiefe) auf derselben Linie wie im Blatt: Das Layout verschiebt sich nicht,
%  solange der Text nicht breiter als die Mindestbreite ist.
\newlength{\lsgbreite}
\newcommand{\lsg}[2][5.25cm]{\leavevmode\ifloesung\settowidth{\lsgbreite}{\,#2\,}\ifdim\lsgbreite<#1\setlength{\lsgbreite}{#1}\fi\rlap{\smash{\raisebox{0.8pt}{\makebox[\lsgbreite][l]{\lsgfarbe\,#2}}}}\luecke[\lsgbreite]\else\luecke[#1]\fi}
% Lücke bis zum Zeilenende (beendet den Absatz): \lsgrest{Lösungstext} bzw. \restzeile
\newcommand{\lsgrest}[1]{\leavevmode\unskip\hspace{1.5mm}\ifloesung\rlap{\smash{\raisebox{0.8pt}{\lsgfarbe\,#1}}}\fi\leaders\hbox{\rule[-0.5ex]{0.5pt}{0.4pt}}\hfill\null\par}
\newcommand{\restzeile}{\lsgrest{}}
% Schreiblinien mit Lösung: \lsglinien{Anzahl}{Lösungstext} (gleiche Höhe wie \linien:
%  je Linie \baselineskip + 5,5 mm + \parskip, ausgemessen)
\newlength{\lsglinienhoehe}
\newcommand{\lsglinien}[2]{\ifloesung\par\vspace{1.5mm}\setlength{\lsglinienhoehe}{\dimexpr #1\dimexpr\baselineskip+5.5mm+\parskip\relax-16.25pt\relax}\noindent\begin{minipage}[t][\lsglinienhoehe][t]{\linewidth}\lsgfarbe #2\end{minipage}\par\vspace{-1.5mm}\else\linien{#1}\fi}
% Beliebiger Inhalt: \lsgoder{nur in der Lösung}{nur im Blatt}
\newcommand{\lsgoder}[2]{\ifloesung\leavevmode#1\else#2\fi}
% Text nur in der Lösung sichtbar, belegt im Blatt aber denselben Platz
\newcommand{\lsgtext}[1]{\leavevmode\ifloesung{\lsgfarbe#1}\else\phantom{#1}\fi}
% Tabellenzelle zum Hineinschreiben (8 mm hoch, ohne Schreiblinie; Lösung mit \lsgtext{…})
\newcommand{\schreibzeile}{\rule[-2.8mm]{0pt}{8mm}}
% Ankreuzfeld: \kreuz{1} = angekreuzt (nur in der Lösung), \kreuz{0} = leer
\newcommand{\kreuz}[1]{\leavevmode\ifloesung\ifnum#1=1 {\lsgfarbe$\boxtimes$}\else$\square$\fi\else$\square$\fi}

% ---------- Syntaxdiagramme (Maße in mm) ----------
%  Terminal rund (orange), Nichtterminal eckig (grün), Linien mit Pfeilen:
%  \draw[sdl,->] (0,0) -- ++(6,0) node[nt,anchor=west] (a) {Augen};   (Pfeil endet am Kasten)
\tikzset{
  sd/.style={x=1mm,y=1mm,>={Stealth[length=1.6mm,width=1.3mm]},font=\small},
  sdl/.style={line width=0.55pt,rounded corners=1.2mm},
  te/.style={draw,rounded corners=2.1mm,minimum height=4.2mm,minimum width=4.2mm,inner sep=1pt,font=\ttfamily\small,fill=orange!18},
  nt/.style={draw,sharp corners,minimum height=4.2mm,inner sep=1.6pt,font=\small,fill=green!12},   % sharp corners: sonst erbt der Kasten die runden Ecken der Linie (sdl)
  sdname/.style={anchor=south west,font=\small\bfseries,inner sep=0pt},
  punkte/.style={font=\small,inner sep=0pt},
  % Ableitungsbaum
  baum/.style={font=\small,edge from parent/.style={draw,line width=0.5pt,-{Stealth[length=1.4mm]}},level distance=8.5mm},
}
% Lösungszeichnung: im Blatt unsichtbar, belegt aber denselben Platz
\ifloesung
  \tikzset{lsgbild/.style={red!75!black,te/.append style={fill=red!5},nt/.append style={fill=red!5}}}
\else
  \tikzset{lsgbild/.style={draw opacity=0,fill opacity=0,text opacity=0}}
\fi
% Karofläche, in der Lösung mit roter Zeichnung: \karolsg{Anzahl Karozeilen}{TikZ-Code}
% Koordinaten der Zeichnung in mm, Ursprung links oben, y nach unten negativ.
\newcommand{\karolsg}[2]{\par\vspace{1mm}\noindent\begin{tikzpicture}
  \draw[step=5mm,gray!45,thin] (0,0) grid (\linewidth-1pt,#1*5mm);
  \useasboundingbox (0,0) rectangle (\linewidth-1pt,#1*5mm);
  \ifloesung\begin{scope}[sd,shift={(0,#1*5mm)},lsgbild]#2\end{scope}\fi
\end{tikzpicture}\par}

% ---------- Bausteine für Syntaxdiagramme (innerhalb tikzpicture[sd]) ----------
%  Jeder Baustein beginnt an einem Punkt (Knoten-Anker oder Koordinate) und
%  legt die Koordinate <name>-o an, an der es weitergeht.
% Aneinanderreihung: \sdnext{name}{von}{te|nt}{Text}   (weiter bei name.east)
\newcommand{\sdnext}[4]{\draw[sdl,->] (#2) -- ++(4,0) node[#3,anchor=west] (#1) {#4};\coordinate (#1-o) at (#1.east);}
% Option: \sdopt{name}{von}{te|nt}{Text}
\newcommand{\sdopt}[4]{\path (#2) ++(4,-5.5) node[#3,anchor=west] (#1) {#4};
  \coordinate (#1-x) at ($(#1.east)+(3.5,0)$); \coordinate (#1-o) at (#1-x |- #2);
  \draw[sdl,->] (#2) |- (#1.west); \draw[sdl,->] (#1.east) -| (#1-o); \draw[sdl] (#2) -- (#1-o);}
% Wiederholung, mindestens einmal: \sdwdh{name}{von}{te|nt}{Text}
\newcommand{\sdwdh}[4]{\draw[sdl,->] (#2) -- ++(6,0) node[#3,anchor=west] (#1) {#4};
  \coordinate (#1-o) at ($(#1.east)+(4,0)$);
  \draw[sdl] (#1.east) -- (#1-o); \draw[sdl,->] (#1-o) -- ++(0,-5.5) -| ($(#2)+(2.5,0)$);}
% Wiederholung, beliebig oft (auch keinmal): \sdwdhn{name}{von}{te|nt}{Text}
\newcommand{\sdwdhn}[4]{\path (#2) ++(6,-5.5) node[#3,anchor=west] (#1) {#4};
  \coordinate (#1-x) at ($(#1.east)+(4,0)$); \coordinate (#1-o) at (#1-x |- #2);
  \draw[sdl] (#2) -- (#1-o); \draw[sdl,->] (#1-o) |- (#1.east); \draw[sdl,->] (#1.west) -| ($(#2)+(2.5,0)$);}
% Alternative: \sdalt{name}{von}{Breite bis zum Zusammenführen}{Zeilenabstand}{te/a, nt/Ziffer, pk/, ...}
%  pk/ setzt senkrechte Pünktchen (…) statt eines Zweigs
\newcommand{\sdalt}[5]{\coordinate (#1-o) at ($(#2)+(#3,0)$);
  \foreach \typ/\txt [count=\i from 0] in {#5}{%
    \ifdefstring{\typ}{pk}{\node[punkte] at ($(#2)+(6.1,-\i*#4+0.4)$) {$\vdots$};}{%
      \ifdefstring{\typ}{te}{\tikzset{sdakt/.style={te}}}{\tikzset{sdakt/.style={nt}}}%
      \node[sdakt,anchor=west] (#1\i) at ($(#2)+(4,-\i*#4)$) {\txt};
      \ifnum\i=0\relax \draw[sdl,->] (#2) -- (#1\i.west); \draw[sdl] (#1\i.east) -- (#1-o);
      \else \draw[sdl,->] (#2) |- (#1\i.west); \draw[sdl,->] (#1\i.east) -| (#1-o);\fi}}}
% Ende eines Diagramms
\newcommand{\sdende}[1]{\draw[sdl,->] (#1) -- ++(5,0);}

% ---------- Die Smiley-Grammatik (Blatt 1.3 und 1.4) ----------
\newcommand{\sdSmileyfolge}{\begin{tikzpicture}[sd,baseline=0mm]
  \node[sdname] at (0,3.3) {Smileyfolge}; \coordinate (s) at (0,0);
  \sdwdh{sm}{s}{nt}{Smiley}\sdende{sm-o}\end{tikzpicture}}
\newcommand{\sdSmiley}{\begin{tikzpicture}[sd,baseline=0mm]
  \node[sdname] at (0,3.3) {Smiley}; \coordinate (s) at (0,0);
  \sdnext{au}{s}{nt}{Augen}\sdnext{mi}{au-o}{nt}{Mitte}\sdnext{mu}{mi-o}{nt}{Mund}\sdende{mu-o}\end{tikzpicture}}
\newcommand{\sdAugen}{\begin{tikzpicture}[sd,baseline=0mm]
  \node[sdname] at (0,3.3) {Augen}; \coordinate (s) at (0,0);
  \sdalt{a}{s}{12}{5.5}{te/:, te/;}\sdende{a-o}\end{tikzpicture}}
\newcommand{\sdMitte}{\begin{tikzpicture}[sd,baseline=0mm]
  \node[sdname] at (0,3.3) {Mitte}; \coordinate (s) at (0,0);
  \sdopt{n}{s}{te}{-}\sdende{n-o}\end{tikzpicture}}
\newcommand{\sdMund}{\begin{tikzpicture}[sd,baseline=0mm]
  \node[sdname] at (0,3.3) {Mund}; \coordinate (s) at (0,0);
  \sdalt{m}{s}{12}{5.5}{te/), te/{(}, te/D, te/P, te/*}\sdende{m-o}\end{tikzpicture}}

% ---------- Ergänzungen für Programmier-Blätter (ab 1.6; aus informatik/10/java/arbeitsblaetter/ab-vorlage.tex) ----------
%  Nur angehängt: Die Makros oben bleiben unverändert, damit die PDFs von 1.1 bis 1.5 gleich bleiben.
%  \lueckentext   am Anfang eines Absatzes/Kastens mit Lücken: Zeilenabstand 1,5, Flattersatz
%  \nr{1}         Schrittnummer ① wie auf der Website und im Code
%  java           Listing:  \begin{java}[numbers=left] … \end{java}
%  javaluecke     Lückencode mit größerem Zeilenabstand, Lücken als (*\lsg[6.75cm]{Lösung}*)
\newcommand{\lueckentext}{\linespread{1.5}\selectfont\raggedright}
\newcommand{\nr}[1]{\ding{\numexpr171+#1\relax}}
\usepackage{listings}
\lstset{language=Java,basicstyle=\ttfamily\small,keywordstyle=\bfseries,commentstyle=\color{gray!80!black},stringstyle=\color{black},showstringspaces=false,columns=flexible,keepspaces=true,tabsize=3,frame=single,framerule=0.4pt,rulecolor=\color{black},xleftmargin=2mm,framexleftmargin=1mm,aboveskip=1.5mm,belowskip=1.5mm,breaklines=true,numberstyle=\tiny\color{gray},numbersep=3mm,escapeinside={(*}{*)},
 literate={ä}{{\"a}}1 {ö}{{\"o}}1 {ü}{{\"u}}1 {Ä}{{\"A}}1 {Ö}{{\"O}}1 {Ü}{{\"U}}1 {ß}{{\ss}}1 {–}{{--}}1 {…}{{\dots}}1}
\lstnewenvironment{java}[1][]{\lstset{#1}}{}
\lstnewenvironment{javaluecke}[1][]{\lstset{basicstyle=\linespread{1.5}\small\ttfamily,#1}}{}

% Blatt:   tectonic ab-1-4-ebnf.tex
% Lösung:  tectonic ab-1-4-ebnf-lsg.tex
% Progression: EBNF-Zeichen aus dem Vergleich mit den bekannten Smiley-Diagrammen erschließen
%              -> Ableitung als Folge von Regelanwendungen -> Diagramm in EBNF übersetzen (Euro)
%              -> Regeln lesen, prüfen, ändern (Molkerei) -> selbst bauen (Uhrzeit)
%              -> Rekursion in der EBNF: Terme aus 1.1 A. 5 (schließt den Bogen „L ⊂ A*, aber wie festlegen?“)
\newcommand{\zeile}{\rule[-1.8mm]{0pt}{6.4mm}}
\newcommand{\ab}[1]{$\stackrel{\mbox{\scriptsize #1}}{\Longrightarrow}$}
\begin{document}
\abtitel{1.4}{Die EBNF}

Die \textbf{EBNF} (erweiterte Backus-Naur-Form) schreibt Produktionsregeln als Text, hier die Smileys aus 1.3:

\vspace{2mm}
\noindent\sdSmileyfolge\hfill\sdSmiley\hfill\sdAugen\hfill\sdMitte\hfill\sdMund\par\vspace{-15mm}
\hspace*{0.5mm}\begin{tabular}{@{}l@{\ \ }l@{}}
R1 & \code{Smileyfolge = Smiley \{Smiley\}.}\\
R2 & \code{Smiley = Augen Mitte Mund.}\\
R3 & \code{Augen = ":" | ";".}\\
R4 & \code{Mitte = ["-"].}\\
R5 & \code{Mund = ")" | "(" | "D" | "P" | "*".}
\end{tabular}
\vspace{1mm}


\begin{merke}[Die EBNF \normalfont(vergleiche jede Regel mit ihrem Diagramm und ergänze)]
\small\renewcommand{\arraystretch}{1.05}
\begin{tabularx}{\linewidth}{@{}p{2.1cm}X|p{1.9cm}X@{}}
\code{Name = … .}\zeile & \lsgoder{\lsgfarbe Produktionsregel für ein Nichtterminal; \code{.} = Ende}{} & \code{"D"}\zeile & \lsgoder{\lsgfarbe Terminal (Zeichen des Alphabets)}{}\\\hline
\code{Mund}\zeile & \lsgoder{\lsgfarbe Nichtterminal (Name einer Regel)}{} & \code{A B}\zeile & \lsgoder{\lsgfarbe Aneinanderreihung: erst A, dann B}{}\\\hline
\code{A | B}\zeile & \lsgoder{\lsgfarbe Alternative: A oder B}{} & \code{[A]}\zeile & \lsgoder{\lsgfarbe Option: A einmal oder keinmal}{}\\\hline
\code{\{A\}}\zeile & \lsgoder{\lsgfarbe Wiederholung: A 0-mal oder öfter}{} & \code{(A | B) C}\zeile & Klammern: \code{("a"|"b")"c"} $\to$ \lsg[1.8cm]{\code{ac}, \code{bc}}\\
\end{tabularx}
\end{merke}

\begin{aufgabe}[Zeichen der EBNF]{1}
\textbf{a)} Warum heißt R1 \code{Smiley \{Smiley\}} und nicht nur \code{\{Smiley\}}?\\[0.6mm]
\lsg[\linewidth]{\code{\{\dots\}} erlaubt auch \emph{keinmal}: Dann wäre \eps{} eine Smileyfolge. Verlangt ist mindestens ein Smiley.}\\[1mm]
\textbf{b)} Schreibe R4 ohne eckige Klammern: \code{Mitte =} \lsg[4.8cm]{\eps\ \code{| "-".}}
\end{aufgabe}

\begin{aufgabe}[Ableitung]{2}
Ableitung von \code{;D:-(}: Je Schritt wird \textbf{ein} Nichtterminal nach seiner Produktionsregel ersetzt.\\[0.5mm]
\begin{minipage}[t]{0.43\linewidth}
\vspace{0pt}
\renewcommand{\arraystretch}{0.92}
\begin{tabular}{@{}l@{\ }l@{}}
& \code{Smileyfolge}\\
\ab{R1} & \code{Smiley Smiley}\\
\ab{R2} & \code{Augen Mitte Mund Smiley}\\
\ab{R3} & \code{";" Mitte Mund Smiley}\\
\ab{R4} & \code{";" Mund Smiley}\qquad {\small (Mitte $\to$ \eps)}\\
\ab{R5} & \code{";" "D" Smiley}\\
\end{tabular}
\end{minipage}\hfill
\begin{minipage}[t]{0.55\linewidth}
\vspace{0pt}
\vspace{2mm}\renewcommand{\arraystretch}{1.3}
\begin{tabularx}{\linewidth}{@{}l@{\ }X@{}}
\ab{\lsg[0.9cm]{R2}} & \lsg[\linewidth]{\code{";" "D" Augen Mitte Mund}}\\
\ab{\lsg[0.9cm]{R3}} & \lsg[\linewidth]{\code{";" "D" ":" Mitte Mund}}\\
\ab{\lsg[0.9cm]{R4}} & \lsg[\linewidth]{\code{";" "D" ":" "-" Mund}}\\
\ab{\lsg[0.9cm]{R5}} & \lsg[\linewidth]{\code{";" "D" ":" "-" "("}}\\
\end{tabularx}
\end{minipage}
\end{aufgabe}

\begin{aufgabe}[Vom Diagramm zur EBNF: Geldbeträge]{3}
Gib die Grammatik der Beträge (\code{12€}, \code{0,79€}, nicht \code{02,99€}) in EBNF an. \code{ZifferOhneNull} ist gegeben.\\[0.5mm]
\begin{tikzpicture}[sd]
  % Geldbetrag
  \node[sdname] at (0,3.3) {Geldbetrag}; \coordinate (g) at (0,0);
  \sdnext{g1}{g}{nt}{Zahl}
  \path (g1-o) ++(4,-5.5) node[te,anchor=west] (g2) {,};
  \draw[sdl,->] (g2.east) -- ++(4,0) node[nt,anchor=west] (g3) {Nachkommastellen};
  \coordinate (g4x) at ($(g3.east)+(3.5,0)$); \coordinate (g4) at (g4x |- g1-o);
  \draw[sdl,->] (g1-o) |- (g2.west); \draw[sdl,->] (g3.east) -| (g4); \draw[sdl] (g1-o) -- (g4);
  \sdnext{g5}{g4}{te}{€}\sdende{g5-o}
  % Zahl
  \node[sdname] at (0,-11.2) {Zahl}; \coordinate (z) at (0,-14.5);
  \draw[sdl,->] (z) -- ++(4,0) node[te,anchor=west] (z0) {0};
  \path (z) ++(4,-5.5) node[nt,anchor=west] (z1) {ZifferOhneNull};
  \draw[sdl,->] (z) |- (z1.west); \coordinate (z1-o) at (z1.east);
  \sdwdhn{z2}{z1-o}{nt}{Ziffer}
  \coordinate (zjx) at ($(z2-o)+(3.5,0)$); \coordinate (zj) at (zjx |- z);
  \draw[sdl] (z0.east) -- (zj); \draw[sdl,->] (z2-o) -| (zj); \sdende{zj}
  % Nachkommastellen
  \node[sdname] at (62,-11.2) {Nachkommastellen}; \coordinate (n) at (62,-14.5);
  \sdnext{n1}{n}{nt}{Ziffer}\sdnext{n2}{n1-o}{nt}{Ziffer}\sdende{n2-o}
  % Ziffer
  \node[sdname] at (120,3.3) {Ziffer}; \coordinate (d) at (120,0);
  \sdalt{d}{d}{31}{5.5}{te/0, nt/ZifferOhneNull}\sdende{d-o}
\end{tikzpicture}\\[-1.5mm]
{\small\renewcommand{\arraystretch}{1.4}
\begin{tabularx}{\linewidth}{@{}r@{\ }X@{}}
\code{Geldbetrag =} & \lsg[\linewidth]{\code{Zahl ["," Nachkommastellen] "€".}}\\
\code{Zahl =} & \lsg[\linewidth]{\code{"0" | ZifferOhneNull \{Ziffer\}.}}\\
\end{tabularx}\\[-0.5mm]
\begin{tabularx}{\linewidth}{@{}r@{\ }X@{\quad}r@{\ }X@{}}
\code{Nachkommastellen =} & \lsg[\linewidth]{\code{Ziffer Ziffer.}} & \code{Ziffer =} & \lsg[\linewidth]{\code{"0" | ZifferOhneNull.}}\\
\end{tabularx}}
\end{aufgabe}

\newpage
\begin{aufgabe}[Regeln lesen und ändern: Molkereinummern]{4}
Ein Code auf Milchprodukten gibt die Molkerei an (Leerzeichen nur zur Lesbarkeit):\\[0.5mm]
\begin{tabular}{@{}l@{\ \ }l@{}}
R1 & \code{Molkereinummer = "DE" Bundesland Betriebsnummer [Europa].}\\
R2 & \code{Bundesland = "BB" | "BE" | "BW" | "BY" | … | "HH" | … | "RP" | … | "TH".}\ \ {\small(16 Länder)}\\
R3 & \code{Betriebsnummer = Ziffer Ziffer Ziffer | Ziffer Ziffer Ziffer Ziffer Ziffer.}\\
R4 & \code{Europa = "EG" | "EU".}\qquad R5\ \ \code{Ziffer = "0" | "1" | … | "9".}
\end{tabular}\\[1.5mm]
\textbf{a)} Gültig? Begründe bei ungültigen Codes.\\[0.8mm]
\code{DE BY 13049 EG} \lsg[3.6cm]{gültig}\hfill \code{DE HH 2514 EU} \lsg[6.75cm]{nein: 4 Ziffern}\\[1.2mm]
\code{BY 70123} \lsg[6.3cm]{nein: \code{DE} fehlt}\\[1mm]
\textbf{b)} Leite \code{DE RP 927} schrittweise ab.
\lsglinien{2}{\code{Molkereinummer} \ab{R1} \code{"DE" Bundesland Betriebsnummer} \ab{R2} \code{"DE" "RP" Betriebsnummer}\\ \ab{R3} \code{"DE" "RP" Ziffer Ziffer Ziffer} \ab{3$\times$R5} \code{"DE" "RP" "9" "2" "7"}\ \ (ohne \code{Europa})}
\textbf{c)} In Bayern gibt die erste Ziffer der Betriebsnummer den Regierungsbezirk an: nur \code{1} bis \code{7} (1 Oberbayern, 2 Niederbayern, …). Ändere und ergänze die Regeln.
\lsglinien{4}{R1 \code{Molkereinummer = "DE" (Bundesland Betriebsnummer | "BY" BetriebsnummerBY) [Europa].}\\
R2 wie bisher, aber \emph{ohne} \code{"BY"}.\\
R6 \code{BetriebsnummerBY = Bezirk Ziffer Ziffer | Bezirk Ziffer Ziffer Ziffer Ziffer.}\\
R7 \code{Bezirk = "1" | "2" | "3" | "4" | "5" | "6" | "7".}}
\end{aufgabe}

\begin{aufgabe}[Selbst bauen: Uhrzeiten]{5}
Uhrzeiten haben die Form \code{hh:mm} oder \code{hh:mm:ss}, immer zweistellig, von \code{00:00} bis \code{23:59:59}.\\[1mm]
\textbf{a)} Drei ungültige Uhrzeiten, die knapp an einer Grenze scheitern: \lsgrest{\code{24:00}, \code{19:60}, \code{12:30:75}}\vspace{0.5mm}
\textbf{b)} Gib eine Grammatik in EBNF an. \textit{Tipp:} Welche Ziffern sind an welcher Stelle erlaubt? Die Stunden ab 20 sind ein Sonderfall.
\lsglinien{4}{\code{Uhrzeit = Stunde ":" Minute [":" Minute].}\quad (Sekunden sind wie Minuten aufgebaut)\\
\code{Stunde = ("0" | "1") Ziffer | "2" ("0" | "1" | "2" | "3").}\\
\code{Minute = ("0" | "1" | "2" | "3" | "4" | "5") Ziffer.}\\
\code{Ziffer = "0" | "1" | "2" | "3" | "4" | "5" | "6" | "7" | "8" | "9".}}
\end{aufgabe}

\begin{aufgabe}[Rekursion in der EBNF: die Terme aus 1.1]{6}
In 1.1 (Aufgabe 5) war \code{sin cos} kein gültiger Term. Diese Grammatik legt die Terme fest:\\[0.5mm]
\begin{tabular}{@{}l@{\ \ }l@{\qquad}l@{\ \ }l@{}}
R1 & \code{Term = Summand \{"+" Summand\}.} & R3 & \code{Funktion = "sin" | "cos" | "sqrt".}\\
R2 & \code{Summand = "x" | Funktion "(" Term ")".} & &
\end{tabular}\\[1mm]
\textbf{a)} Gültig? Begründe bei ungültigen Wörtern.\\[0.6mm]
\code{sqrt(x+x)} \lsg[2.25cm]{ja}\hfill \code{cos(sin(x))+x} \lsg[2.25cm]{ja}\hfill \code{sin(x)+} \lsg[5.1cm]{nein: Summand fehlt}\\[1.2mm]
\code{(x+x)} \lsg[5.4cm]{nein: \code{(} ohne Funktion}\\[1mm]
\textbf{b)} Welche Regeln verweisen über einen Umweg auf sich selbst? \lsgrest{R1, R2: \code{Term} → \code{Summand} → \code{Term}}\vspace{0.5mm}
\textbf{c)} (mündlich) Warum kommt man hier ohne Rekursion nicht aus? Vergleiche mit 1.3, Aufgabe 5d.\\[1mm]
\textbf{d)} Ändere R2, sodass auch \code{(x+x)} gültig ist:\\[1mm]
\code{Summand =} \lsg[10.8cm]{\code{"x" | [Funktion] "(" Term ")".}}
\end{aufgabe}

\end{document}
